
\Data\ID{1003154401}
\Updated{2020-07-15 03:07:21}
\Level{A}
\SubArea{Power and radical functions}
\LANGen{
\Question{
Find the true statement about the function \( f(x)=2-(x-1)^3 \).}
{The function \( f \) is an injective (one-to-one) function.}{The function \( f \) is increasing.}{The function \( f \) is odd.}{The function \( f \) has the maximum at \( x=1 \).}{}{}}
\LANGpl{
\Question{
Które stwierdzenie dotyczące funkcji  \( f(x)=2-(x-1)^3 \) jest prawdziwe?}
{Funkcja  \( f \) jest iniekcyjna (jeden do jednego).}{Funkcja  \( f \) jest rosnąca.}{Funkcja \( f \) jest nieparzysta.}{Funkcja  \( f \) osiąga maksimum w  \( x=1 \).}{}{}}
\LANGcs{
\Question{
Funkce \( f \) je dána předpisem \( f(x)=2-(x-1)^3 \). Vyberte pravdivý výrok.}
{Funkce \( f \) je prostá.}{Funkce \( f \) je rostoucí.}{Funkce \( f \) je lichá.}{Funkce \( f \) má maximum v bodě \( x=1 \).}{}{}}
\LANGsk{
\Question{
Vyberte pravdivý výrok o funkcií \( f(x)=2-(x-1)^3 \).}
{Funkcia \( f \) je prostá.}{Funkcia \( f \) je rastúca.}{Funkcia \( f \) je nepárna.}{Funkcia \( f \) má maximum v bode \( x=1 \).}{}{}}
\LANGes{
\Question{
Determina la proposición lógica verdadera sobre la función \( f(x)=2-(x-1)^3 \).}
{La función \( f \) es una función inyectiva (uno a uno).}{La función  \( f \) es creciente.}{La función \( f \) es impar.}{La función  \( f \) tiene el máximo en \( x=1 \).}{}{}}
\End

\Data\ID{1003154402}
\Updated{2020-07-15 03:07:21}
\Level{A}
\SubArea{Power and radical functions}
\LANGen{
\Question{
Find the false statement about the function \( f(x)=3-(x+2)^4 \).}
{The function \( f \) is even.}{The function \( f \) has the maximum at \( x=-2 \).}{The function \( f \) is bounded above.}{The range of the function \( f \) is the interval \( (-\infty; 3] \).}{}{}}
\LANGpl{
\Question{
Które z poniższych stwierdzeń dotyczących funkcji \( f(x)=3-(x+2)^4 \) jest fałszywe?}
{Funkcja \( f \) is parzysta.}{Funkcja \( f \) osiąga maksimum w \( x=-2 \).}{Funkcja \( f \) jest ograniczona z góry.}{Zakres funkcji \( f \) mieści się w przedziale \( (-\infty; 3\rangle \).}{}{}}
\LANGcs{
\Question{
Funkce \( f \) je dána předpisem \( f(x)=3-(x+2)^4 \). Vyberte nepravdivý výrok.}
{Funkce \( f \) je sudá.}{Funkce \( f \) má maximum v bodě \( x=-2 \).}{Funkce \( f \) je shora omezená.}{Oborem hodnot funkce \( f \) je interval \( (-\infty; 3\rangle \).}{}{}}
\LANGsk{
\Question{
Vyberte nepravdivý výrok o funkcií \( f(x)=3-(x+2)^4 \).}
{Funkcia \( f \) je párna.}{Funkcia \( f \) má maximum v bode \( x=-2 \).}{Funkcia \( f \) je zhora ohraničená.}{Obor hodnôt funkcie \( f \) je interval \( (-\infty; 3\rangle \).}{}{}}
\LANGes{
\Question{
Determina la proposición lógica falsa sobre la función \( f(x)=3-(x+2)^4 \).}
{La función  \( f \) es par.}{La función \( f \) tiene el máximo en  \( x=-2 \).}{La función \( f \) está acotada superiormente.}{El rango de la función  \( f \) es el intervalo  \( (-\infty; 3] \).}{}{}}
\End

\Data\ID{1003162201}
\Updated{2020-07-15 03:07:21}
\Level{A}
\SubArea{Power and radical functions}
\LANGen{
\Question{
Let \( f(x)=x^{-2} \). Identify which of the following statements is false.}
{\( f\left(\frac23\right)=\frac49 \)}{\( f(-0.125)=64 \)}{\( f\left(\frac14\right)=16 \)}{\( f\left(\frac1{0.5}\right)=\frac14 \)}{}{}}
\LANGpl{
\Question{
Niech \( f(x)=x^{-2} \). Które z poniższych stwierdzeń jest fałszywe?}
{\( f\left(\frac23\right)=\frac49 \)}{\( f(-0{,}125)=64 \)}{\( f\left(\frac14\right)=16 \)}{\( f\left(\frac1{0{,}5}\right)=\frac14 \)}{}{}}
\LANGcs{
\Question{
Funkce \(f\) je dána předpisem \( f(x)=x^{-2} \). Vyberte nepravdivý výrok.}
{\( f\left(\frac23\right)=\frac49 \)}{\( f(-0{,}125)=64 \)}{\( f\left(\frac14\right)=16 \)}{\( f\left(\frac1{0{,}5}\right)=\frac14 \)}{}{}}
\LANGsk{
\Question{
Funkcia \(f\) je daná predpisom \( f(x)=x^{-2} \). Vyberte nepravdivý výrok.}
{\( f\left(\frac23\right)=\frac49 \)}{\( f(-0{,}125)=64 \)}{\( f\left(\frac14\right)=16 \)}{\( f\left(\frac1{0{,}5}\right)=\frac14 \)}{}{}}
\LANGes{
\Question{
Sea \( f(x)=x^{-2} \). Identifica cuál de las siguientes proposiciones lógicas es falsa.}
{\( f\left(\frac23\right)=\frac49 \)}{\( f(-0.125)=64 \)}{\( f\left(\frac14\right)=16 \)}{\( f\left(\frac1{0.5}\right)=\frac14 \)}{}{}}
\End

\Data\ID{1003162202}
\Updated{2020-07-15 03:07:21}
\Level{A}
\SubArea{Power and radical functions}
\LANGen{
\Question{
Let \( f(x)=x^{-3} \). Identify which of the following statements is true.}
{\( f\left(-\frac23\right)=-\frac{27}8 \)}{\( f(0)=1 \)}{\( f\left(-\frac1{10}\right)=-0.001 \)}{\( f(0.2)=\frac1{125} \)}{}{}}
\LANGpl{
\Question{
Niech  \( f(x)=x^{-3} \). Które z poniższych stwierdzeń jest prawdziwe?}
{\( f\left(-\frac23\right)=-\frac{27}8 \)}{\( f(0)=1 \)}{\( f\left(-\frac1{10}\right)=-0{,}001 \)}{\( f(0{,}2)=\frac1{125} \)}{}{}}
\LANGcs{
\Question{
Funkce \(f\) je dána předpisem \( f(x)=x^{-3} \). Vyberte pravdivý výrok.}
{\( f\left(-\frac23\right)=-\frac{27}8 \)}{\( f(0)=1 \)}{\( f\left(-\frac1{10}\right)=-0{,}001 \)}{\( f(0{,}2)=\frac1{125} \)}{}{}}
\LANGsk{
\Question{
Funkcia \(f\) je daná predpisom \( f(x)=x^{-3} \). Vyberte pravdivý výrok.}
{\( f\left(-\frac23\right)=-\frac{27}8 \)}{\( f(0)=1 \)}{\( f\left(-\frac1{10}\right)=-0{,}001 \)}{\( f(0{,}2)=\frac1{125} \)}{}{}}
\LANGes{
\Question{
Sea \( f(x)=x^{-3} \). Identifica cuál de las siguientes proposiciones lógicas es verdadera.}
{\( f\left(-\frac23\right)=-\frac{27}8 \)}{\( f(0)=1 \)}{\( f\left(-\frac1{10}\right)=-0.001 \)}{\( f(0.2)=\frac1{125} \)}{}{}}
\End

\Data\ID{1003162203}
\Updated{2020-07-15 03:07:21}
\Level{A}
\SubArea{Power and radical functions}
\LANGen{
\Question{
Let \( f(x)=x^{-2} \)  and  \( g(x)=x^{-3} \). Identify which of the following statements is false.}
{\( f(2)+g(2)=\frac1{32} \)}{\( f(2)\cdot g(2)=\frac1{32} \)}{\( \frac{f(2)}{g(2)} =2 \)}{\( f\!\left(2^{-3}\right)=64 \)}{}{}}
\LANGpl{
\Question{
Niech  \( f(x)=x^{-2} \)  i \( g(x)=x^{-3} \). Które z poniższych stwierdzeń jest fałszywe?}
{\( f(2)+g(2)=\frac1{32} \)}{\( f(2)\cdot g(2)=\frac1{32} \)}{\( \frac{f(2)}{g(2)} =2 \)}{\( f\!\left(2^{-3}\right)=64 \)}{}{}}
\LANGcs{
\Question{
Funkce \(f\) a \(g\) jsou dány předpisy \( f(x)=x^{-2} \); \( g(x)=x^{-3} \). Vyberte nepravdivý výrok.}
{\( f(2)+g(2)=\frac1{32} \)}{\( f(2)\cdot g(2)=\frac1{32} \)}{\( \frac{f(2)}{g(2)} =2 \)}{\( f\!\left(2^{-3}\right)=64 \)}{}{}}
\LANGsk{
\Question{
Funkcie \(f\) a \(g\) sú dané predpismi \( f(x)=x^{-2} \) a \( g(x)=x^{-3} \). Vyberte nepravdivý výrok.}
{\( f(2)+g(2)=\frac1{32} \)}{\( f(2)\cdot g(2)=\frac1{32} \)}{\( \frac{f(2)}{g(2)} =2 \)}{\( f\!\left(2^{-3}\right)=64 \)}{}{}}
\LANGes{
\Question{
Sean \( f(x)=x^{-2} \)  y \( g(x)=x^{-3} \). Identifica cuál de las siguientes proposiciones lógicas es falsa.}
{\( f(2)+g(2)=\frac1{32} \)}{\( f(2)\cdot g(2)=\frac1{32} \)}{\( \frac{f(2)}{g(2)} =2 \)}{\( f\!\left(2^{-3}\right)=64 \)}{}{}}
\End

\Data\ID{1103143401}
\Updated{2021-04-07 09:04:49}
\Level{A}
\SubArea{Power and radical functions}
\IMG{1103143401.pdf}
\LANGen{
\Question{
The graphs represent the parts of the functions \( f(x)=x^3 \)  and \( g(x)=x^4 \). Identify which of the following statements is false.}
{\( \left(-\frac12\right)^3 > (2)^3 \)}{\( 	(-2)^3 < \left(-\frac12\right)^3 \)}{\( \left(\frac13\right)^3 \geq (0.3)^3 \)}{\( 	(-1)^3 \leq (1)^3 \)}{}{}}
\LANGpl{
\Question{
Wykresy reprezentują części funkcji \( f(x)=x^3 \)  i \( g(x)=x^4 \). Które z poniższych stwierdzeń jest fałszywe?}
{\( \left(-\frac12\right)^3 > (2)^3 \)}{\( 	(-2)^3 < \left(-\frac12\right)^3 \)}{\( \left(\frac13\right)^3 \geq (0{,}3)^3 \)}{\( 	(-1)^3 \leq (1)^3 \)}{}{}}
\LANGcs{
\Question{
Na obrázku jsou části grafů funkcí \( f(x)=x^3 \) a \( g(x)=x^4 \). Vyberte nepravdivý výrok.}
{\( \left(-\frac12\right)^3 > (2)^3 \)}{\( 	(-2)^3 < \left(-\frac12\right)^3 \)}{\( \left(\frac13\right)^3 \geq (0{,}3)^3 \)}{\( 	(-1)^3 \leq (1)^3 \)}{}{}}
\LANGsk{
\Question{
Na obrázku sú časti grafov funkcií \( f(x)=x^3 \)  a \( g(x)=x^4 \). Vyberte nepravdivý výrok.}
{\( \left(-\frac12\right)^3 > (2)^3 \)}{\( 	(-2)^3 < \left(-\frac12\right)^3 \)}{\( \left(\frac13\right)^3 \geq (0{,}3)^3 \)}{\( 	(-1)^3 \leq (1)^3 \)}{}{}}
\LANGes{
\Question{
Las gráficas representan las partes de las funciones \( f(x)=x^3 \)  y \( g(x)=x^4 \). Identifica cuál de las siguientes proposiciones lógicas es  falsa.}
{\( \left(-\frac12\right)^3 > (2)^3 \)}{\( 	(-2)^3 < \left(-\frac12\right)^3 \)}{\( \left(\frac13\right)^3 \geq (0.3)^3 \)}{\( 	(-1)^3 \leq (1)^3 \)}{}{}}
\End

\Data\ID{1103143402}
\Updated{2021-04-07 09:04:27}
\Level{A}
\SubArea{Power and radical functions}
\IMG{1103143402.pdf}
\LANGen{
\Question{
The graphs represent the parts of the functions \( f(x)=x^3 \) and \( g(x)=x^4 \). Identify which of the following statements is true.}
{\( (-3)^4 > (2)^4 \)}{\( (-2)^4 < \left(-\frac12\right)^4 \)}{\( \left(-\frac14\right)^4 \geq (0.3)^4 \)}{\( (-1)^4 < (1)^4 \)}{}{}}
\LANGpl{
\Question{
Wykresy reprezentują części funkcji \( f(x)=x^3 \) i \( g(x)=x^4 \). Które z poniższych stwierdzeń jest prawdziwe?}
{\( (-3)^4 > (2)^4 \)}{\( (-2)^4 < \left(-\frac12\right)^4 \)}{\( \left(-\frac14\right)^4 \geq (0{,}3)^4 \)}{\( (-1)^4 < (1)^4 \)}{}{}}
\LANGcs{
\Question{
Na obrázku jsou části grafů funkcí \( f(x)=x^3 \) a \( g(x)=x^4 \). Vyberte pravdivý výrok.}
{\( (-3)^4 > (2)^4 \)}{\( (-2)^4 < \left(-\frac12\right)^4 \)}{\( \left(-\frac14\right)^4 \geq (0{,}3)^4 \)}{\( (-1)^4 < (1)^4 \)}{}{}}
\LANGsk{
\Question{
Na obrázku sú časti grafov funkcií \( f(x)=x^3 \) a \( g(x)=x^4 \). Vyberte pravdivý výrok.}
{\( (-3)^4 > (2)^4 \)}{\( (-2)^4 < \left(-\frac12\right)^4 \)}{\( \left(-\frac14\right)^4 \geq (0{,}3)^4 \)}{\( (-1)^4 < (1)^4 \)}{}{}}
\LANGes{
\Question{
Las gráficas representan las partes de las funciones \( f(x)=x^3 \) y \( g(x)=x^4 \). Identifica cuál de las siguientes proposiciones lógicas es verdadera.}
{\( (-3)^4 > (2)^4 \)}{\( (-2)^4 < \left(-\frac12\right)^4 \)}{\( \left(-\frac14\right)^4 \geq (0.3)^4 \)}{\( (-1)^4 < (1)^4 \)}{}{}}
\End

\Data\ID{1103143403}
\Updated{2021-04-07 09:04:46}
\Level{A}
\SubArea{Power and radical functions}
\IMG{1103143403.pdf}
\LANGen{
\Question{
The graphs represent the parts of the functions \( f(x)=x^3 ;\ g(x)=x^4;\ h(x)=x^5 \). Identify which of the following statements is false.}
{\( \left(-\frac13\right)^5 < \left(-\frac13\right)^3 \)}{\( \left(\frac12\right)^5 < \left(-\frac12\right)^4 \)}{\( (-3)^4 > (3)^3 \)}{\( \left(\frac14\right)^3 \geq (-0.25)^4 \)}{}{}}
\LANGpl{
\Question{
Poniższe wykresy reprezentują części funkcji \( f(x)=x^3 ;\ g(x)=x^4;\ h(x)=x^5 \). Które z poniższych stwierdzeń jest fałszywe?}
{\( \left(-\frac13\right)^5 < \left(-\frac13\right)^3 \)}{\( \left(\frac12\right)^5 < \left(-\frac12\right)^4 \)}{\( (-3)^4 > (3)^3 \)}{\( \left(\frac14\right)^3 \geq (-0{,}25)^4 \)}{}{}}
\LANGcs{
\Question{
Na obrázku jsou části grafů funkcí  \( f(x)=x^3 \); \( g(x)=x^4;\ h(x)=x^5 \). Vyberte nepravdivý výrok.}
{\( \left(-\frac13\right)^5 < \left(-\frac13\right)^3 \)}{\( \left(\frac12\right)^5 < \left(-\frac12\right)^4 \)}{\( (-3)^4 > (3)^3 \)}{\( \left(\frac14\right)^3 \geq (-0{,}25)^4 \)}{}{}}
\LANGsk{
\Question{
Na obrázku sú časti grafov funkcií \( f(x)=x^3 ;\ g(x)=x^4;\ h(x)=x^5 \). Vyberte nepravdivý výrok.}
{\( \left(-\frac13\right)^5 < \left(-\frac13\right)^3 \)}{\( \left(\frac12\right)^5 < \left(-\frac12\right)^4 \)}{\( (-3)^4 > (3)^3 \)}{\( \left(\frac14\right)^3 \geq (-0{,}25)^4 \)}{}{}}
\LANGes{
\Question{
Las gráficas representan las partes de las funciones \( f(x)=x^3 ;\ g(x)=x^4;\ h(x)=x^5 \). Identifica cuál de las siguientes proposiciones lógicas es falsa.}
{\( \left(-\frac13\right)^5 < \left(-\frac13\right)^3 \)}{\( \left(\frac12\right)^5 < \left(-\frac12\right)^4 \)}{\( (-3)^4 > (3)^3 \)}{\( \left(\frac14\right)^3 \geq (-0.25)^4 \)}{}{}}
\End

\Data\ID{1103143501}
\Updated{2021-04-07 09:04:00}
\Level{A}
\SubArea{Power and radical functions}
\IMG{1103143501.pdf}
\LANGen{
\Question{
The graphs represent the parts of the functions \( f(x)=x^3 \) and \( g(x)=x^4 \). Identify which of the following statements is false.}
{The solution set of the inequality \( x^4 > x^3 \) is \( (1;\infty) \).}{The solution set of the inequality \( x^4 > 0 \) is \( (-\infty;0)\cup(0;\infty) \).}{The solution set of the equation \( x^3 = x^4 \) is \( \{0;1\} \).}{The solution set of the inequality \( x^3 \geq x^4 \) is \( [0;1] \).}{}{}}
\LANGpl{
\Question{
Poniższe wykresy reprezentują części funkcji \( f(x)=x^3 \) i \( g(x)=x^4 \). Które z poniższych stwierdzeń jest fałszywe?}
{Zbiorem rozwiązań nierówności \( x^4 > x^3 \) jest  \( (1;\infty) \).}{Zbiorem rozwiązań nierówności \( x^4 > 0 \) jest  \( (-\infty;0)\cup(0;\infty) \).}{Zbiorem rozwiązań równania  \( x^3 = x^4 \) jest \( \{0;1\} \).}{Zbiorem rozwiązań nierówności \( x^3 \geq x^4 \) jest  \( \langle0;1\rangle \).}{}{}}
\LANGcs{
\Question{
Na obrázku jsou části grafů funkcí \( f(x)=x^3 \) a \( g(x)=x^4 \). Vyberte nepravdivý výrok.}
{Množinou všech řešení nerovnice \( x^4 > x^3 \) je \( (1;\infty) \).}{Množinou všech řešení nerovnice \( x^4 > 0 \) je \( (-\infty;0)\cup(0;\infty) \).}{Množinou všech řešení rovnice \( x^3 = x^4 \) je \( \{0;1\} \).}{Množinou všech řešení nerovnice \( x^3 \geq x^4 \) je \( \langle0;1\rangle \).}{}{}}
\LANGsk{
\Question{
Na obrázku sú časti grafov funkcií \( f(x)=x^3 \) a \( g(x)=x^4 \). Vyberte nepravdivý výrok.}
{Množina všetkých riešení nerovnice \( x^4 > x^3 \) je \( (1;\infty) \).}{Množina všetkých riešení nerovnice \( x^4 > 0 \) je \( (-\infty;0)\cup(0;\infty) \).}{Množina všetkých riešení rovnice \( x^3 = x^4 \) je \( \{0;1\} \).}{Množina všetkých riešení nerovnice \( x^3 \geq x^4 \) je \( \langle0;1\rangle \).}{}{}}
\LANGes{
\Question{
Las gráficas representan las partes de las funciones  \( f(x)=x^3 \) y \( g(x)=x^4 \). Identifica cuál de las siguientes proposiciones lógicas es falsa.}
{El conjunto de soluciones de la desigualdad \( x^4 > x^3 \) es \( (1;\infty) \).}{El conjunto de soluciones de la desigualdad \( x^4 > 0 \) es \( (-\infty;0)\cup(0;\infty) \).}{El conjunto de soluciones de la ecuación \( x^3 = x^4 \) es \( \{0;1\} \).}{El conjunto de soluciones de la desigualdad  \( x^3 \geq x^4 \) es \( [0;1] \).}{}{}}
\End

\Data\ID{1103143502}
\Updated{2021-04-07 09:04:17}
\Level{A}
\SubArea{Power and radical functions}
\IMG{1103143502.pdf}
\LANGen{
\Question{
The graphs represent the parts of the functions \( f(x)=x^3 \) and \( g(x)=x^5 \). Identify which of the following inequalities  has the solution set \( (-1;0)\cup(1;\infty) \).}
{\( x^3 < x^5 \)}{\( x^5 \geq x^3 \)}{\( x^3 > x^5 \)}{\( x^3 > -1 \)}{}{}}
\LANGpl{
\Question{
Wykresy reprezentują części funkcji \( f(x)=x^3 \) i \( g(x)=x^5 \). Zbiorem rozwiązań której z poniższych nierówności jest  \( (-1;0)\cup(1;\infty) \)?}
{\( x^3 < x^5 \)}{\( x^5 \geq x^3 \)}{\( x^3 > x^5 \)}{\( x^3 > -1 \)}{}{}}
\LANGcs{
\Question{
Na obrázku jsou části grafů funkcí \( f(x)=x^3 \) a \( g(x)=x^5 \). Vyberte nerovnici, jejíž množinou všech řešení je \( (-1;0)\cup(1;\infty) \).}
{\( x^3 < x^5 \)}{\( x^5 \geq x^3 \)}{\( x^3 > x^5 \)}{\( x^3 > -1 \)}{}{}}
\LANGsk{
\Question{
Na obrázku sú časti grafov funkcií \( f(x)=x^3 \) a \( g(x)=x^5 \). Určte, ktorá z nasledujúcich nerovníc má množinu všetkých riešení \( (-1;0)\cup(1;\infty) \).}
{\( x^3 < x^5 \)}{\( x^5 \geq x^3 \)}{\( x^3 > x^5 \)}{\( x^3 > -1 \)}{}{}}
\LANGes{
\Question{
Las gráficas representan las partes de las funciones  \( f(x)=x^3 \) y \( g(x)=x^5 \). Identifica cuál de las siguientes desigualdades tiene el conjunto de soluciones \( (-1;0)\cup(1;\infty) \).}
{\( x^3 < x^5 \)}{\( x^5 \geq x^3 \)}{\( x^3 > x^5 \)}{\( x^3 > -1 \)}{}{}}
\End

\Data\ID{1103143503}
\Updated{2021-04-07 09:04:31}
\Level{A}
\SubArea{Power and radical functions}
\IMG{1103143503.pdf}
\LANGen{
\Question{
The graphs represent the parts of the functions \( f(x)=x^4 \)  and \( g(x)=x^6 \). Identify which of the following statements is true.}
{The solution set of the inequality \( x^4 \leq x^6 \) is \( (-\infty; -1]\cup[1;\infty)\cup\{0\} \).}{The solution set of the inequality \( x^4 > x^6 \) is \( (-1;1) \).}{The solution set of the equation \( x^6=x^4 \) is  \( \{0;1\} \).}{The solution set of the inequality \( x^6 \geq x^4 \) is \( (-\infty; -1]\cup[1; \infty) \).}{}{}}
\LANGpl{
\Question{
Wykresy reprezentują części funkcji \( f(x)=x^4 \)  i \( g(x)=x^6 \). Które z poniższych stwierdzeń jest prawdziwe?}
{Zbiorem rozwiązań nierówności   \( x^4 \leq x^6 \) jest \( (-\infty; -1\rangle\cup\langle1;\infty)\cup\{0\} \).}{Zbiorem rozwiązań nierówności  \( x^4 > x^6 \) jest \( (-1;1) \).}{Zbiorem rozwiązań równania  \( x^6=x^4 \) jest \( \{0;1\} \).}{Zbiorem rozwiązań nierówności  \( x^6 \geq x^4 \) jest \( (-\infty; -1\rangle\cup\langle1; \infty) \).}{}{}}
\LANGcs{
\Question{
Na obrázku jsou části grafů funkcí \( f(x)=x^4 \)  a \( g(x)=x^6 \). Vyberte pravdivý výrok.}
{Množinou všech řešení nerovnice \( x^4 \leq x^6 \) je \( (-\infty; -1\rangle\cup\langle1;\infty)\cup\{0\} \).}{Množinou všech řešení nerovnice \( x^4 > x^6 \) je \( (-1;1) \).}{Množinou všech řešení rovnice \( x^6=x^4 \) je \( \{0;1\} \).}{Množinou všech řešení nerovnice \( x^6 \geq x^4 \) je \( (-\infty; -1\rangle\cup\langle1; \infty) \).}{}{}}
\LANGsk{
\Question{
Na obrázku sú časti grafov funkcií  \( f(x)=x^4 \)  a \( g(x)=x^6 \). Vyberte pravdivý výrok.}
{Množina všetkých riešení nerovnice \( x^4 \leq x^6 \) je \( (-\infty; -1\rangle\cup\langle1;\infty)\cup\{0\} \).}{Množina všetkých riešení nerovnice \( x^4 > x^6 \) je \( (-1;1) \).}{Množina všetkých riešení rovnice \( x^6=x^4 \) je  \( \{0;1\} \).}{Množina všetkých riešení nerovnice \( x^6 \geq x^4 \) je \( (-\infty; -1\rangle\cup\langle1; \infty) \).}{}{}}
\LANGes{
\Question{
Las gráficas representan las partes de las funciones \( f(x)=x^4 \)  y \( g(x)=x^6 \). Identifica cuál de las siguientes proposiciones lógicas es verdadera.}
{El conjunto de soluciones de la desigualdad  \( x^4 \leq x^6 \) es \( (-\infty; -1]\cup[1;\infty)\cup\{0\} \).}{El conjunto de soluciones de la desigualdad \( x^4 > x^6 \) es \( (-1;1) \).}{El conjunto de soluciones de la ecuación \( x^6=x^4 \) es  \( \{0;1\} \).}{El conjunto de soluciones de la desigualdad  \( x^6 \geq x^4 \) es \( (-\infty; -1]\cup[1; \infty) \).}{}{}}
\End

\Data\ID{1103156801}
\Updated{2020-07-15 03:07:21}
\Level{A}
\SubArea{Power and radical functions}
\IMG{1103156801.pdf}
\LANGen{
\Question{
Function \( f \) is given by the graph. Identify the equation of the function \( f \).}
{\( f(x)=8-(x-1)^3;\ x\in[0;3] \)}{\( f(x)=9-x^5;\ x\in[0;3] \)}{\( f(x)=\left|(x+1)^2+8\right|;\ x\in[0;3] \)}{\( f(x)=(x+1)^3+8;\ x\in[0;3] \)}{}{}}
\LANGpl{
\Question{
Funkcja \( f \) przedstawiona została za pomocą wykresu. Wyznacz równanie funkcji \( f \).}
{\( f(x)=8-(x-1)^3;\ x\in\langle0;3\rangle \)}{\( f(x)=9-x^5;\ x\in\langle0;3\rangle \)}{\( f(x)=\left|(x+1)^2+8\right|;\ x\in\langle0;3\rangle \)}{\( f(x)=(x+1)^3+8;\ x\in\langle0;3\rangle \)}{}{}}
\LANGcs{
\Question{
Funkce \( f \) je dána grafem. Vyberte předpis funkce \( f \).}
{\( f(x)=8-(x-1)^3;\ x\in\langle0;3\rangle \)}{\( f(x)=9-x^5;\ x\in\langle0;3\rangle \)}{\( f(x)=\left|(x+1)^2+8\right|;\ x\in\langle0;3\rangle \)}{\( f(x)=(x+1)^3+8;\ x\in\langle0;3\rangle \)}{}{}}
\LANGsk{
\Question{
Funkcia \( f \) je daná grafom. Vyberte predpis funkcie \( f \).}
{\( f(x)=8-(x-1)^3;\ x\in\langle0;3\rangle \)}{\( f(x)=9-x^5;\ x\in\langle0;3\rangle \)}{\( f(x)=\left|(x+1)^2+8\right|;\ x\in\langle0;3\rangle \)}{\( f(x)=(x+1)^3+8;\ x\in\langle0;3\rangle \)}{}{}}
\LANGes{
\Question{
La función  \( f \) viene dada por el gráfico. Identifica la ecuación de la función \( f \).}
{\( f(x)=8-(x-1)^3;\ x\in[0;3] \)}{\( f(x)=9-x^5;\ x\in[0;3] \)}{\( f(x)=\left|(x+1)^2+8\right|;\ x\in[0;3] \)}{\( f(x)=(x+1)^3+8;\ x\in[0;3] \)}{}{}}
\End

\Data\ID{1103156802}
\Updated{2020-07-15 03:07:21}
\Level{A}
\SubArea{Power and radical functions}
\IMG{1103156802.pdf}
\LANGen{
\Question{
The graphs represent functions \( f(x)=x^3 \), \( x\in[0;1]\); \( g(x)=2x^3 \), \( x\in[0;1] \) and \( h(x)=x^4 \), \( x\in[0;1] \). 
Choose the legend which assigns the correct color of the graph to each of the given functions.}
{\( f \) -- red, \( g \) --  green, \( h \) -- blue}{\( f \) --  blue, \( g \) --  green, \( h \) -- red}{\( f \) -- blue, \( g \) --  red, \( h \) -- green}{\( f \) --  green, \( g \) --  red, \( h \) -- blue}{}{}}
\LANGpl{
\Question{
Wykresy reprezentują funkcje \( f(x)=x^3 \), \( x\in\langle0;1\rangle\); \( g(x)=2x^3 \), \( x\in\langle0;1\rangle \) i \( h(x)=x^4 \), \( x\in\langle0;1\rangle \). 
Wybierz legendę, która w poprawny sposób przyporządkowuje wykres do każdej z podanych funkcji.}
{\( f \) -- czerwony, \( g \) --  zielony, \( h \) -- niebieski}{\( f \) --  niebieski, \( g \) --  zielony, \( h \) -- czerwony}{\( f \) -- niebieski, \( g \) --  czerwony, \( h \) -- zielony}{\( f \) --  zielony, \( g \) --  czerwony, \( h \) -- niebieski}{}{}}
\LANGcs{
\Question{
Na obrázku jsou grafy tří funkcí, které jsou dány předpisy: \( f(x)=x^3 \), \( x\in\langle0;1\rangle\); \( g(x)=2x^3 \), \( x\in\langle0;1\rangle \); \( h(x)=x^4 \), \( x\in\langle0;1\rangle \). 
Vyberte odpověď, v níž jsou grafům funkcí přiřazeny správné barvy.}
{\( f \) - červená, \( g \) - zelená, \( h \) - modrá}{\( f \) - modrá, \( g \) - zelená, \( h \) - červená}{\( f \) - modrá, \( g \) - červená, \( h \) - zelená}{\( f \) - zelená, \( g \) - červená, \( h \) - modrá}{}{}}
\LANGsk{
\Question{
Na obrázku sú grafy troch funkcií \( f(x)=x^3 \), \( x\in\langle0;1\rangle\); \( g(x)=2x^3 \), \( x\in\langle0;1\rangle \) a \( h(x)=x^4 \), \( x\in\langle0;1\rangle \). 
Vyberte možnosť, v ktorej sú grafom daných funkcií priradené správne farby.}
{\( f \) -- červená, \( g \) --  zelená, \( h \) -- modrá}{\( f \) --  modrá, \( g \) --  zelená, \( h \) -- červená}{\( f \) -- modrá, \( g \) --  červená, \( h \) -- zelená}{\( f \) --  zelená, \( g \) --  červená, \( h \) -- modrá}{}{}}
\LANGes{
\Question{
Las gráficas representan funciones \( f(x)=x^3 \), \( x\in[0;1]\); \( g(x)=2x^3 \), \( x\in[0;1] \) y \( h(x)=x^4 \), \( x\in[0;1] \). 
Elige la leyenda que asigna el color correcto de la gráfica de cada una de las funciones dadas.}
{\( f \) -- rojo, \( g \) --  verde, \( h \) -- azul}{\( f \) --  azul, \( g \) --  verde, \( h \) -- rojo}{\( f \) -- azul, \( g \) --  rojo, \( h \) -- verde}{\( f \) --  verde, \( g \) --  rojo, \( h \) -- azul}{}{}}
\End

\Data\ID{1103156803}
\Updated{2020-07-15 03:07:21}
\Level{A}
\SubArea{Power and radical functions}
\IMG{1103156803.pdf}
\LANGen{
\Question{
The graphs represent the functions  \( f(x)=x^2 \), \( x\in[-2;0] \); \( g(x)=x^3\), \(x\in[-2;0] \)  and \( h(x)=(-x)^3 \), \( x\in[-2;0] \). 
Choose the legend which assigns the correct color of the graph to each of the given functions.}
{\( f \) --  blue, \( g \) -- red, \( h \) -- green}{\( f \) --  red, \( g \) --  blue, \( h \) -- green}{\( f \) --  green, \( g \) --  red, \( h \) -- blue}{\( f \) --  blue, \( g \) --  green, \( h \) -- red}{}{}}
\LANGpl{
\Question{
Wykresy reprezentują funkcje  \( f(x)=x^2 \), \( x\in\langle-2;0\rangle \); \( g(x)=x^3\), \(x\in\langle-2;0\rangle \)  i \( h(x)=(-x)^3 \), \( x\in\langle-2;0\rangle \). 
Wybierz legendę, która w poprawny sposób przyporządkowuje kolor wykresu do funkcji.}
{\( f \) --  niebieski, \( g \) -- czerwony, \( h \) -- zielony}{\( f \) --  czerwony, \( g \) --  niebieski, \( h \) -- zielony}{\( f \) --  zielony, \( g \) --  czerwony, \( h \) -- niebieski}{\( f \) --  niebieski, \( g \) --  zielony, \( h \) -- czerwony}{}{}}
\LANGcs{
\Question{
Na obrázku jsou grafy tří funkcí, které jsou dány předpisy: \( f(x)=x^2 \), \( x\in\langle-2;0\rangle \); \( g(x)=x^3\), \(x\in\langle-2;0\rangle \); \( h(x)=(-x)^3 \), \( x\in\langle-2;0\rangle \). 
Vyberte odpověď, v níž jsou grafům funkcí přiřazeny správné barvy.}
{\( f \) - modrá, \( g \) - červená, \( h \) - zelená}{\( f \) - červená, \( g \) - modrá, \( h \) - zelená}{\( f \) - zelená, \( g \) - červená, \( h \) - modrá}{\( f \) - modrá, \( g \) - zelená, \( h \) - červená}{}{}}
\LANGsk{
\Question{
Na obrázku sú grafy troch funkcií \( f(x)=x^2 \), \( x\in\langle-2;0\rangle \); \( g(x)=x^3\), \(x\in\langle-2;0\rangle \)  a \( h(x)=(-x)^3 \), \( x\in\langle-2;0\rangle \). 
Vyberte možnosť, v ktorej sú grafom daných funkcií priradené správne farby.}
{\( f \) --  modrá, \( g \) -- červená, \( h \) -- zelená}{\( f \) --  červená, \( g \) --  modrá, \( h \) -- zelená}{\( f \) --  zelená, \( g \) --  červená, \( h \) -- modrá}{\( f \) --  modrá, \( g \) --  zelená, \( h \) -- červená}{}{}}
\LANGes{
\Question{
Las gráficas representan funciones\( f(x)=x^2 \), \( x\in[-2;0] \); \( g(x)=x^3\), \(x\in[-2;0] \)  y \( h(x)=(-x)^3 \), \( x\in[-2;0] \). 
Elige la leyenda que asigna el color correcto de la gráfica de cada una de las funciones dadas.}
{\( f \) --  azul, \( g \) -- rojo, \( h \) -- verde}{\( f \) --  rojo, \( g \) --  azul, \( h \) -- verde}{\( f \) --  verde, \( g \) --  rojo, \( h \) -- azul}{\( f \) --  azul, \( g \) --  verde, \( h \) -- rojo}{}{}}
\End

\Data\ID{1103159301}
\Updated{2021-07-24 10:07:30}
\Level{A}
\SubArea{Power and radical functions}
\IMG{1103159301_0.pdf}
\LANGen{
\Question{
The graphs represent the parts of the functions \( f(x)=x^{-2} \) and \( g(x)=x^{-3} \). Identify which of the following statements is false.}
{\( \left(\frac12\right)^{-3} < 2^{-3} \)}{\( \left(-\frac12\right)^{-3} < 2^{-3} \)}{\( \left( -\frac12\right)^{-2} \geq (-2)^{-2} \)}{\( (-2)^{-2} \geq 2^{-2} \)}{}{}}
\LANGpl{
\Question{
Wykresy poniżej przedstawiają części funkcji  \( f(x)=x^{-2} \) and \( g(x)=x^{-3} \). Które z poniższych stwierdzeń jest  fałszywe?}
{\( \left(\frac12\right)^{-3} < 2^{-3} \)}{\( \left(-\frac12\right)^{-3} < 2^{-3} \)}{\( \left( -\frac12\right)^{-2} \geq (-2)^{-2} \)}{\( (-2)^{-2} \geq 2^{-2} \)}{}{}}
\LANGcs{
\Question{
Na obrázku jsou části grafů funkcí \( f(x)=x^{-2} \) a \( g(x)=x^{-3} \). Vyberte nepravdivý výrok.}
{\( \left(\frac12\right)^{-3} < 2^{-3} \)}{\( \left(-\frac12\right)^{-3} < 2^{-3} \)}{\( \left( -\frac12\right)^{-2} \geq (-2)^{-2} \)}{\( (-2)^{-2} \geq 2^{-2} \)}{}{}}
\LANGsk{
\Question{
Na obrázku sú časti grafov funkcií \( f(x)=x^{-2} \) a \( g(x)=x^{-3} \). Vyberte nepravdivý výrok.}
{\( \left(\frac12\right)^{-3} < 2^{-3} \)}{\( \left(-\frac12\right)^{-3} < 2^{-3} \)}{\( \left( -\frac12\right)^{-2} \geq (-2)^{-2} \)}{\( (-2)^{-2} \geq 2^{-2} \)}{}{}}
\LANGes{
\Question{
Las gráficas representan las partes de las funciones \( f(x)=x^{-2} \) y \( g(x)=x^{-3} \). Identifica cuál de las siguientes proposiciones lógicas es  falsa.}
{\( \left(\frac12\right)^{-3} < 2^{-3} \)}{\( \left(-\frac12\right)^{-3} < 2^{-3} \)}{\( \left( -\frac12\right)^{-2} \geq (-2)^{-2} \)}{\( (-2)^{-2} \geq 2^{-2} \)}{}{}}
\End

\Data\ID{1103159302}
\Updated{2020-07-15 03:07:21}
\Level{A}
\SubArea{Power and radical functions}
\IMG{1103159302.pdf}
\LANGen{
\Question{
The graphs represent the parts of the functions \( f(x)=x^{-3} \)  and \( g(x)=x^{-4} \). Identify which of the following statements is true.}
{\( \left(\frac12\right)^{-3} < \left( \frac12 \right)^{-4} \)}{\( 2^{-4} > 2^{-3} \)}{\( (-2)^{-4} \leq (-2)^{-3} \)}{\( (-1)^{-4} > 1^{-3} \)}{}{}}
\LANGpl{
\Question{
Wykresy reprezentują części funkcji \( f(x)=x^{-3} \)  i \( g(x)=x^{-4} \). Które z poniższych stwierdzeń jest prawdziwe?}
{\( \left(\frac12\right)^{-3} < \left( \frac12 \right)^{-4} \)}{\( 2^{-4} > 2^{-3} \)}{\( (-2)^{-4} \leq (-2)^{-3} \)}{\( (-1)^{-4} > 1^{-3} \)}{}{}}
\LANGcs{
\Question{
Na obrázku jsou části grafů funkcí \( f(x)=x^{-3} \)  a \( g(x)=x^{-4} \). Vyberte pravdivý výrok.}
{\( \left(\frac12\right)^{-3} < \left( \frac12 \right)^{-4} \)}{\( 2^{-4} > 2^{-3} \)}{\( (-2)^{-4} \leq (-2)^{-3} \)}{\( (-1)^{-4} > 1^{-3} \)}{}{}}
\LANGsk{
\Question{
Na obrázku sú časti grafov funkcií \( f(x)=x^{-3} \)  a \( g(x)=x^{-4} \). Vyberte pravdivý výrok.}
{\( \left(\frac12\right)^{-3} < \left( \frac12 \right)^{-4} \)}{\( 2^{-4} > 2^{-3} \)}{\( (-2)^{-4} \leq (-2)^{-3} \)}{\( (-1)^{-4} > 1^{-3} \)}{}{}}
\LANGes{
\Question{
Las gráficas representan las partes de las funciones \( f(x)=x^{-3} \) y \( g(x)=x^{-4} \). Identifica cuál de las siguientes proposiciones lógicas es verdadera.}
{\( \left(\frac12\right)^{-3} < \left( \frac12 \right)^{-4} \)}{\( 2^{-4} > 2^{-3} \)}{\( (-2)^{-4} \leq (-2)^{-3} \)}{\( (-1)^{-4} > 1^{-3} \)}{}{}}
\End

\Data\ID{1103159303}
\Updated{2020-07-15 03:07:21}
\Level{A}
\SubArea{Power and radical functions}
\IMG{1103159303_0.pdf}
\LANGen{
\Question{
The graphs represent the parts of the functions \( f(x)=x^{-2} \)  and \( g(x)=x^{-3} \). Identify which of the following statements is true.}
{\( -\left(\frac12\right)^{-3} < (-2)^{-3} \)}{\( (-2)^{-2} \leq -2^{-2} \)}{\( (-2)^{-3} < -2^{-3} \)}{\( (-2)^{-3} \leq -2^{-2} \)}{}{}}
\LANGpl{
\Question{
Wykresy reprezentują części funkcji \( f(x)=x^{-2} \)  i \( g(x)=x^{-3} \). Które z poniższych stwierdzeń jest prawdziwe?}
{\( -\left(\frac12\right)^{-3} < (-2)^{-3} \)}{\( (-2)^{-2} \leq -2^{-2} \)}{\( (-2)^{-3} < -2^{-3} \)}{\( (-2)^{-3} \leq -2^{-2} \)}{}{}}
\LANGcs{
\Question{
Na obrázku jsou části grafů funkcí  \( f(x)=x^{-2} \)  a \( g(x)=x^{-3} \). Vyberte pravdivý výrok.}
{\( -\left(\frac12\right)^{-3} < (-2)^{-3} \)}{\( (-2)^{-2} \leq -2^{-2} \)}{\( (-2)^{-3} < -2^{-3} \)}{\( (-2)^{-3} \leq -2^{-2} \)}{}{}}
\LANGsk{
\Question{
Na obrázku sú časti grafov funkcií \( f(x)=x^{-2} \)  a \( g(x)=x^{-3} \). Vyberte pravdivý výrok.}
{\( -\left(\frac12\right)^{-3} < (-2)^{-3} \)}{\( (-2)^{-2} \leq -2^{-2} \)}{\( (-2)^{-3} < -2^{-3} \)}{\( (-2)^{-3} \leq -2^{-2} \)}{}{}}
\LANGes{
\Question{
Las gráficas representan las partes de las funciones \( f(x)=x^{-2} \)  y \( g(x)=x^{-3} \). Identifica cuál de las siguientes proposiciones lógicas es verdadera.}
{\( -\left(\frac12\right)^{-3} < (-2)^{-3} \)}{\( (-2)^{-2} \leq -2^{-2} \)}{\( (-2)^{-3} < -2^{-3} \)}{\( (-2)^{-3} \leq -2^{-2} \)}{}{}}
\End

\Data\ID{1103161001}
\Updated{2020-07-15 03:07:21}
\Level{A}
\SubArea{Power and radical functions}
\IMG{1103161001.pdf}
\LANGen{
\Question{
The graphs represent the parts of the functions \( f(x)=x^{-2} \)  and \( g(x)=x^{-3} \). Identify which of the following statements is false.}
{The solution set of the inequality \( x^{-2} > 0 \) is \( (-\infty;\infty) \).}{The solution set of the inequality \( x^{-3} > 0 \) is \( (0;\infty) \).}{The solution set of the equation \( x^{-3} = x^{-2} \) is \( \{1\} \).}{The solution set of the inequality \( x^{-3} < x^{-2} \) is \( (-\infty;0)\cup(1;\infty) \).}{}{}}
\LANGpl{
\Question{
Wykresy reprezentują części funkcji \( f(x)=x^{-2} \)  i \( g(x)=x^{-3} \). Które z poniższych stwierdzeń jest  fałszywe?}
{Zbiorem rozwiązań nierówności  \( x^{-2} > 0 \) jest  \( (-\infty;\infty) \).}{Zbiorem rozwiązań nierówności  \( x^{-3} > 0 \) jest \( (0;\infty) \).}{Zbiorem rozwiązań równania \( x^{-3} = x^{-2} \) jest \( \{1\} \).}{Zbiorem rozwiązań nierówności \( x^{-3} < x^{-2} \) jest  \( (-\infty;0)\cup(1;\infty) \).}{}{}}
\LANGcs{
\Question{
Na obrázku jsou části grafů funkcí \( f(x)=x^{-2} \)  a \( g(x)=x^{-3} \). Vyberte nepravdivý výrok.}
{Množinou všech řešení nerovnice \( x^{-2} > 0 \) je \( (-\infty;\infty) \).}{Množinou všech řešení nerovnice \( x^{-3} > 0 \) je \( (0;\infty) \).}{Množinou všech řešení rovnice \( x^{-3} = x^{-2} \) je \( \{1\} \).}{Množinou všech řešení nerovnice \( x^{-3} < x^{-2} \) je \( (-\infty;0)\cup(1;\infty) \).}{}{}}
\LANGsk{
\Question{
Na obrázku sú časti grafov funkcií  \( f(x)=x^{-2} \)  a \( g(x)=x^{-3} \). Vyberte nepravdivý výrok.}
{Množina všetkých riešení nerovnice \( x^{-2} > 0 \) je \( (-\infty;\infty) \).}{Množina všetkých riešení nerovnice \( x^{-3} > 0 \) je \( (0;\infty) \).}{Množina všetkých riešení rovnice \( x^{-3} = x^{-2} \) je \( \{1\} \).}{Množina všetkých riešení nerovnice \( x^{-3} < x^{-2} \) je \( (-\infty;0)\cup(1;\infty) \).}{}{}}
\LANGes{
\Question{
Las gráficas representan las partes de las funciones \( f(x)=x^{-2} \)  y \( g(x)=x^{-3} \). Identifica cuál de las siguientes proposiciones lógicas es falsa.}
{El conjunto de soluciones de la desigualdad \( x^{-2} > 0 \) es \( (-\infty;\infty) \).}{El conjunto de soluciones de la desigualdad \( x^{-3} > 0 \) es \( (0;\infty) \).}{El conjunto de soluciones de la ecuación\( x^{-3} = x^{-2} \) es \( \{1\} \).}{El conjunto de soluciones de la desigualdad \( x^{-3} < x^{-2} \) es \( (-\infty;0)\cup(1;\infty) \).}{}{}}
\End

\Data\ID{1103161002}
\Updated{2021-04-07 09:04:27}
\Level{A}
\SubArea{Power and radical functions}
\IMG{1103161002.pdf}
\LANGen{
\Question{
The graphs represent the parts of the functions \( f(x)=x^{-2} \)  and \( g(x)=x^{-4} \). Identify which of the following inequalities has the solution set \( (-\infty; -1]\cup[1;\infty) \).}
{\( x^{-4} \leq x^{-2} \)}{\( x^{-2} \leq x^{-4} \)}{\( x^{-2} > x^{-4} \)}{\( x^{-2} < 1 \)}{}{}}
\LANGpl{
\Question{
Wykresy reprezentują części funkcji \( f(x)=x^{-2} \)  i \( g(x)=x^{-4} \). Zbiorem rozwiązań której z poniższych nierówności jest \( (-\infty; -1\rangle\cup\langle1;\infty) \).}
{\( x^{-4} \leq x^{-2} \)}{\( x^{-2} \leq x^{-4} \)}{\( x^{-2} > x^{-4} \)}{\( x^{-2} < 1 \)}{}{}}
\LANGcs{
\Question{
Na obrázku jsou části grafů funkcí \( f(x)=x^{-2} \)  a \( g(x)=x^{-4} \). Vyberte nerovnici, jejíž množinou všech řešení je \( (-\infty; -1\rangle\cup\langle1;\infty) \).}
{\( x^{-4} \leq x^{-2} \)}{\( x^{-2} \leq x^{-4} \)}{\( x^{-2} > x^{-4} \)}{\( x^{-2} < 1 \)}{}{}}
\LANGsk{
\Question{
Na obrázku sú časti grafov funkcií \( f(x)=x^{-2} \)  a \( g(x)=x^{-4} \). Určte, ktorá z nasledujúcich nerovníc má množinu všetkých riešení \( (-\infty; -1\rangle\cup\langle1;\infty) \).}
{\( x^{-4} \leq x^{-2} \)}{\( x^{-2} \leq x^{-4} \)}{\( x^{-2} > x^{-4} \)}{\( x^{-2} < 1 \)}{}{}}
\LANGes{
\Question{
Las gráficas representan las partes de las funciones \( f(x)=x^{-2} \)  y \( g(x)=x^{-4} \). Identifica cuál de las siguientes desigualdades tiene como conjunto de soluciones  \( (-\infty; -1]\cup[1;\infty) \).}
{\( x^{-4} \leq x^{-2} \)}{\( x^{-2} \leq x^{-4} \)}{\( x^{-2} > x^{-4} \)}{\( x^{-2} < 1 \)}{}{}}
\End

\Data\ID{1103161003}
\Updated{2020-07-15 03:07:21}
\Level{A}
\SubArea{Power and radical functions}
\IMG{1103161003.pdf}
\LANGen{
\Question{
The graphs represent the parts of the functions \( f(x)=x^{-2} \)  and \( g(x)=x^{-3} \). Identify which of the following inequalities has the solution set \( (-\infty;-1)\cup(0;\infty) \).}
{\( -x^{-3} < x^{-2} \)}{\( \left|x^{-3}\right| < x^{-2} \)}{\( x^{-3} < -x^{-2} \)}{\( x^{-3} < \left|x^{-2}\right| \)}{}{}}
\LANGpl{
\Question{
Wykresy reprezentują części funkcji \( f(x)=x^{-2} \)  i \( g(x)=x^{-3} \). Zbiorem rozwiązań której z podanych nierówności jest \( (-\infty;-1)\cup(0;\infty) \).}
{\( -x^{-3} < x^{-2} \)}{\( \left|x^{-3}\right| < x^{-2} \)}{\( x^{-3} < -x^{-2} \)}{\( x^{-3} < \left|x^{-2}\right| \)}{}{}}
\LANGcs{
\Question{
Na obrázku jsou části grafů funkcí \( f(x)=x^{-2} \)  a \( g(x)=x^{-3} \). Vyberte nerovnici, jejíž množinou všech řešení je \( (-\infty;-1)\cup(0;\infty) \).}
{\( -x^{-3} < x^{-2} \)}{\( \left|x^{-3}\right| < x^{-2} \)}{\( x^{-3} < -x^{-2} \)}{\( x^{-3} < \left|x^{-2}\right| \)}{}{}}
\LANGsk{
\Question{
Na obrázku sú časti grafov funkcií \( f(x)=x^{-2} \)  a \( g(x)=x^{-3} \). Určte, ktorá z nasledujúcich nerovníc má množinu všetkých riešení \( (-\infty;-1)\cup(0;\infty) \).}
{\( -x^{-3} < x^{-2} \)}{\( \left|x^{-3}\right| < x^{-2} \)}{\( x^{-3} < -x^{-2} \)}{\( x^{-3} < \left|x^{-2}\right| \)}{}{}}
\LANGes{
\Question{
Las gráficas representan las partes de las funciones \( f(x)=x^{-2} \)  y  \( g(x)=x^{-3} \). Identifica cuál de las siguientes desigualdades tiene el conjunto de soluciones  \( (-\infty;-1)\cup(0;\infty) \).}
{\( -x^{-3} < x^{-2} \)}{\( \left|x^{-3}\right| < x^{-2} \)}{\( x^{-3} < -x^{-2} \)}{\( x^{-3} < \left|x^{-2}\right| \)}{}{}}
\End

\Data\ID{1103163101}
\Updated{2021-04-07 09:04:04}
\Level{A}
\SubArea{Power and radical functions}
\IMG{1103163101.pdf}
\LANGen{
\Question{
Function \( f \) is given by the graph. Identify which of the following statements is true.}
{\( f(x)=2+(x-1)^{-2};\ x\in[1.5;6] \)}{\( f(x)=2+(x-2)^{-2};\ x\in[1.5;6] \)}{\( f(x)=2+(x-1)^2;\ x\in[1.5;6] \)}{\( f(x)=2+(x-1)^{-1};\ x\in[1.5;6] \)}{}{}}
\LANGpl{
\Question{
Funkcję  \( f \) przedstawiono za pomocą wykresu. Które z poniższych stwierdzeń jest prawdziwe?}
{\( f(x)=2+(x-1)^{-2};\ x\in\langle1{,}5;6\rangle \)}{\( f(x)=2+(x-2)^{-2};\ x\in\langle1{,}5;6\rangle \)}{\( f(x)=2+(x-1)^2;\ x\in\langle1{,}5;6\rangle \)}{\( f(x)=2+(x-1)^{-1};\ x\in\langle1{,}5;6\rangle \)}{}{}}
\LANGcs{
\Question{
Funkce \( f \) je dána grafem. Vyberte pravdivý výrok.}
{\( f(x)=2+(x-1)^{-2};\ x\in\langle1{,}5;6\rangle \)}{\( f(x)=2+(x-2)^{-2};\ x\in\langle1{,}5;6\rangle \)}{\( f(x)=2+(x-1)^2;\ x\in\langle1{,}5;6\rangle \)}{\( f(x)=2+(x-1)^{-1};\ x\in\langle1{,}5;6\rangle \)}{}{}}
\LANGsk{
\Question{
Funkcia \( f \) je daná grafom. Vyberte pravdivý výrok.}
{\( f(x)=2+(x-1)^{-2};\ x\in\langle1{,}5;6\rangle \)}{\( f(x)=2+(x-2)^{-2};\ x\in\langle1{,}5;6\rangle \)}{\( f(x)=2+(x-1)^2;\ x\in\langle1{,}5;6\rangle \)}{\( f(x)=2+(x-1)^{-1};\ x\in\langle1{,}5;6\rangle \)}{}{}}
\LANGes{
\Question{
La función \( f \) viene dada por el gráfico. Identifica cuál de las siguientes proposiciones lógicas es verdadera.}
{\( f(x)=2+(x-1)^{-2};\ x\in[1.5;6] \)}{\( f(x)=2+(x-2)^{-2};\ x\in[1.5;6] \)}{\( f(x)=2+(x-1)^2;\ x\in[1.5;6] \)}{\( f(x)=2+(x-1)^{-1};\ x\in[1.5;6] \)}{}{}}
\End

\Data\ID{1103163102}
\Updated{2020-07-15 03:07:21}
\Level{A}
\SubArea{Power and radical functions}
\IMG{1103163102.pdf}
\LANGen{
\Question{
Function \( f \) is given by the graph. 
Identify which of the following statements is true.}
{\( f(x)=(x+1)^{-3};\ x\in[-3;-1.5] \)}{\( f(x)=-(x+1)^{-2};\ x\in[-3;-1.5] \)}{\( f(x)=(x+1)^{-5};\ x\in[-3;-1.5] \)}{\( f(x)=(x+1)^{-1};\ x\in[-3;-1.5] \)}{}{}}
\LANGpl{
\Question{
Funkcja \( f \) przedstawiona jest za pomocą wykresu.
Które z poniższych stwierdzeń jest prawdziwe?}
{\( f(x)=(x+1)^{-3};\ x\in\langle-3;-1{,}5\rangle \)}{\( f(x)=-(x+1)^{-2};\ x\in\langle-3;-1{,}5\rangle \)}{\( f(x)=(x+1)^{-5};\ x\in\langle-3;-1{,}5\rangle \)}{\( f(x)=(x+1)^{-1};\ x\in\langle-3;-1{,}5\rangle \)}{}{}}
\LANGcs{
\Question{
Funkce \( f \) je dána grafem. 
Vyberte pravdivý výrok.}
{\( f(x)=(x+1)^{-3};\ x\in\langle-3;-1{,}5\rangle \)}{\( f(x)=-(x+1)^{-2};\ x\in\langle-3;-1{,}5\rangle \)}{\( f(x)=(x+1)^{-5};\ x\in\langle-3;-1{,}5\rangle \)}{\( f(x)=(x+1)^{-1};\ x\in\langle-3;-1{,}5\rangle \)}{}{}}
\LANGsk{
\Question{
Funkcia \( f \) je daná grafom. Vyberte pravdivý výrok.}
{\( f(x)=(x+1)^{-3};\ x\in\langle-3;-1{,}5\rangle \)}{\( f(x)=-(x+1)^{-2};\ x\in\langle-3;-1{,}5\rangle \)}{\( f(x)=(x+1)^{-5};\ x\in\langle-3;-1{,}5\rangle \)}{\( f(x)=(x+1)^{-1};\ x\in\langle-3;-1{,}5\rangle \)}{}{}}
\LANGes{
\Question{
La función \( f \) viene dada por el gráfico.
Identifica cuál de las siguientes proposiciones lógicas es verdadera.}
{\( f(x)=(x+1)^{-3};\ x\in[-3;-1.5] \)}{\( f(x)=-(x+1)^{-2};\ x\in[-3;-1.5] \)}{\( f(x)=(x+1)^{-5};\ x\in[-3;-1.5] \)}{\( f(x)=(x+1)^{-1};\ x\in[-3;-1.5] \)}{}{}}
\End

\Data\ID{1103163103}
\Updated{2020-07-15 03:07:21}
\Level{A}
\SubArea{Power and radical functions}
\IMG{1103163103.pdf}
\LANGen{
\Question{
The parts of the graphs represent the functions: \( f(x)=x^{-2} \), \( g(x)=x^{-3} \), \( h(x)=x^{-4} \).
Choose the legend which assigns the correct color of the graph to each of the given functions.}
{\( f \) -- green, \( g \)  --  blue, \( h \) -- red}{\( f \)  --  red, \( g \) --  blue, \( h \) -- green}{\( f \) -- green, \( g \) --  red, \( h \) -- blue}{\( f \) --  blue, \( g \) --  green, \( h \) -- red}{}{}}
\LANGpl{
\Question{
Części wykresów reprezentują funkcje: \( f(x)=x^{-2} \), \( g(x)=x^{-3} \), \( h(x)=x^{-4} \).
Wybierz legendę, która przypisuje odpowiednie kolory wykresu  do każdej z podanych funkcji.}
{\( f \) -- zielony, \( g \)  --  niebieski, \( h \) -- czerwony}{\( f \)  --  czerwony, \( g \) --  niebieski, \( h \) -- zielony}{\( f \) -- zielony, \( g \) --  czerwony, \( h \) -- niebieski}{\( f \) --  niebieski, \( g \) --  zielony, \( h \) -- czerwony}{}{}}
\LANGcs{
\Question{
Na obrázku jsou části grafů tří funkcí, které jsou dány předpisy: \( f(x)=x^{-2} \), \( g(x)=x^{-3} \), \( h(x)=x^{-4} \).
Vyberte odpověď, v níž jsou grafům funkcí přiřazeny správné barvy.}
{\( f \) - zelená, \( g \)  - modrá, \( h \) - červená}{\( f \)  - červená, \( g \) - modrá, \( h \) - zelená}{\( f \) - zelená, \( g \) - červená, \( h \) - modrá}{\( f \) - modrá, \( g \) - zelená, \( h \) - červená}{}{}}
\LANGsk{
\Question{
Na obrázku sú časti grafov troch funkcií \( f(x)=x^{-2} \), \( g(x)=x^{-3} \), \( h(x)=x^{-4} \). Vyberte možnosť, v ktorej sú grafom daných funkcií priradené správne farby.}
{\( f \) -- zelená, \( g \)  --  modrá, \( h \) -- červená}{\( f \)  --  červená, \( g \) --  modrá, \( h \) -- zelená}{\( f \) -- zelená, \( g \) --  červená, \( h \) -- modrá}{\( f \) --  modrá, \( g \) --  zelená, \( h \) -- červená}{}{}}
\LANGes{
\Question{
Las gráficas representan las partes de las funciones: \( f(x)=x^{-2} \), \( g(x)=x^{-3} \), \( h(x)=x^{-4} \).
Elige la leyenda que asigna el color correcto de la gráfica de cada una de las funciones dadas.}
{\( f \) -- verde, \( g \)  --  azul, \( h \) -- rojo}{\( f \)  --  rojo, \( g \) --  azul, \( h \) -- verde}{\( f \) -- verde, \( g \) --  rojo, \( h \) -- azul}{\( f \) --  azul, \( g \) --  verde, \( h \) -- rojo}{}{}}
\End

\Data\ID{9000010605}
\Updated{2021-04-07 09:04:06}
\Level{A}
\SubArea{Power and radical functions}
\LANGen{
\Question{
Identify a function which is decreasing on
\((-\infty ;1)\).}
{\(f(x) = -x^{3}\)}{\(f(x) = -x^{2}\)}{\(f(x) = x^{3}\)}{\(f(x) = -x^{4}\)}{\(f(x)= x^{-2}\)}{\(f(x) = x^{2}\)}}
\LANGpl{
\Question{
Wybierz funkcję, która maleje w przedziale
\((-\infty ;1)\).}
{\(f \colon y= -x^{3}\)}{\(f \colon y = -x^{2}\)}{\(f \colon y = x^{3}\)}{\(f \colon y = -x^{4}\)}{\(f \colon y = x^{-2}\)}{\(f \colon y = x^{2}\)}}
\LANGcs{
\Question{
Vyberte předpis funkce, která je v intervalu
\((-\infty ;1)\) 
klesající.}
{\[f \colon y = -x^{3}\]}{\[f \colon y = -x^{2}\]}{\[f \colon y = x^{3}\]}{\[f \colon y = -x^{4}\]}{\[f \colon y = x^{-2}\]}{\[f \colon y = x^{2}\]}}
\LANGsk{
\Question{
Určte funkciu, ktorá je na intervale
\((-\infty ;1)\) 
klesajúca.}
{\(f \colon y = -x^{3}\)}{\(f \colon y = -x^{2}\)}{\(f \colon y = x^{3}\)}{\(f \colon y = -x^{4}\)}{\(f \colon y = x^{-2}\)}{\(f \colon y = x^{2}\)}}
\LANGes{
\Question{
Identifica una función que está decreciendo en
\((-\infty ;1)\).}
{\(f(x) = -x^{3}\)}{\(f(x) = -x^{2}\)}{\(f(x) = x^{3}\)}{\(f(x) = -x^{4}\)}{\(f(x)= x^{-2}\)}{\(f(x) = x^{2}\)}}
\End

\Data\ID{9000010606}
\Updated{2021-04-07 09:04:32}
\Level{A}
\SubArea{Power and radical functions}
\LANGen{
\Question{
Identify a function which is increasing on
\((-1;3)\).}
{\(y = (x + 2)^{2}\)}{\(y = x^{2} + x\)}{\(y = x^{2} - x\)}{\(y = (x - 2)^{2}\)}{\(y = -x^{3}\)}{\(y = x^{2} + 1\)}}
\LANGpl{
\Question{
Wybierz funkcję rosnącą w przedziale
\((-1;3)\).}
{\(y = (x + 2)^{2}\)}{\(y = x^{2} + x\)}{\(y = x^{2} - x\)}{\(y = (x - 2)^{2}\)}{\(y = -x^{3}\)}{\(y = x^{2} + 1\)}}
\LANGcs{
\Question{
Vyberte funkci, která je v intervalu
\((-1;3)\) 
rostoucí.}
{\(y = (x + 2)^{2}\)}{\(y = x^{2} + x\)}{\(y = x^{2} - x\)}{\(y = (x - 2)^{2}\)}{\(y = -x^{3}\)}{\(y = x^{2} + 1\)}}
\LANGsk{
\Question{
Určte funkciu, ktorá je na intervale 
\((-1;3)\) 
rastúca.}
{\(y = (x + 2)^{2}\)}{\(y = x^{2} + x\)}{\(y = x^{2} - x\)}{\(y = (x - 2)^{2}\)}{\(y = -x^{3}\)}{\(y = x^{2} + 1\)}}
\LANGes{
\Question{
Identifica una función que está creciendo en
\((-1;3)\).}
{\(y = (x + 2)^{2}\)}{\(y = x^{2} + x\)}{\(y = x^{2} - x\)}{\(y = (x - 2)^{2}\)}{\(y = -x^{3}\)}{\(y = x^{2} + 1\)}}
\End

\Data\ID{9000010607}
\Updated{2021-04-07 09:04:51}
\Level{A}
\SubArea{Power and radical functions}
\LANGen{
\Question{
Identify a function which is one-to-one on the interval
\([ - 2;2] \).}
{\(y = x^{3} - 2\)}{\(y = x^{2} - 2\)}{\(y = -x^{2} + 2\)}{\(y = x^{-2} + 2\)}{\(y = \frac{1} 
{x} - 2\)}{\(y = x^{4}\)}}
\LANGpl{
\Question{
Która z funkcji poniżej jest prosta w przedziale
\([ - 2;2] \)?}
{\(y = x^{3} - 2\)}{\(y = x^{2} - 2\)}{\(y = -x^{2} + 2\)}{\(y = x^{-2} + 2\)}{\(y = \frac{1} 
{x} - 2\)}{\(y = x^{4}\)}}
\LANGcs{
\Question{
Vyberte funkci, která je v intervalu
\(\langle - 2;2\rangle \) 
prostá.}
{\(y = x^{3} - 2\)}{\(y = x^{2} - 2\)}{\(y = -x^{2} + 2\)}{\(y = x^{-2} + 2\)}{\(y = \frac{1} 
{x} - 2\)}{\(y = x^{4}\)}}
\LANGsk{
\Question{
Určte funkciu, ktorá je na intervale 
\(\langle - 2;2\rangle \) 
prostá.}
{\(y = x^{3} - 2\)}{\(y = x^{2} - 2\)}{\(y = -x^{2} + 2\)}{\(y = x^{-2} + 2\)}{\(y = \frac{1} 
{x} - 2\)}{\(y = x^{4}\)}}
\LANGes{
\Question{
Identifica una función que sea uno a uno (inyectiva) en el intervalo
\([ - 2;2] \).}
{\(y = x^{3} - 2\)}{\(y = x^{2} - 2\)}{\(y = -x^{2} + 2\)}{\(y = x^{-2} + 2\)}{\(y = \frac{1} 
{x} - 2\)}{\(y = x^{4}\)}}
\End

\Data\ID{9000011103}
\Updated{2020-07-15 03:07:34}
\Level{A}
\SubArea{Power and radical functions}
\LANGen{
\Question{
In the following list identify an increasing function.}
{\(f(x) = x^{5}\)}{\(f(x) = x^{2}\)}{\(f(x) = x^{-3}\)}{\(f(x) = x^{-4}\)}{\(f(x) = 2x^{0}\)}{}}
\LANGpl{
\Question{
Która z poniższych funkcji jest funkcją rosnącą?}
{\(f\colon y = x^{5}\)}{\(f\colon y = x^{2}\)}{\(f\colon y = x^{-3}\)}{\(f\colon y = x^{-4}\)}{\(f\colon y = 2x^{0}\)}{}}
\LANGcs{
\Question{
Z následujících funkcí vyberte tu, která je rostoucí v celém svém
definičním oboru:}
{\(f\colon y = x^{5}\)}{\(f\colon y = x^{2}\)}{\(f\colon y = x^{-3}\)}{\(f\colon y = x^{-4}\)}{\(f\colon y = 2x^{0}\)}{}}
\LANGsk{
\Question{
Ktorá z nasledujúcich funkcií je rastúca na celom svojom definičnom obore?}
{\(f\colon y = x^{5}\)}{\(f\colon y = x^{2}\)}{\(f\colon y = x^{-3}\)}{\(f\colon y = x^{-4}\)}{\(f\colon y = 2x^{0}\)}{}}
\LANGes{
\Question{
Identifica una función creciente en la siguiente lista.}
{\(f(x) = x^{5}\)}{\(f(x) = x^{2}\)}{\(f(x) = x^{-3}\)}{\(f(x) = x^{-4}\)}{\(f(x) = 2x^{0}\)}{}}
\End

\Data\ID{9000025801}
\Updated{2020-07-15 03:07:34}
\Level{A}
\SubArea{Power and radical functions}
\LANGen{
\Question{
Find all intersections of the graph of the following function with
\(x\)-axis.
\[ 
 f(x) = x^{3} - x^{2} - 2x
\]}
{\(X_{1} = [0;0]\),
\(X_{2} = [-1;0]\),
\(X_{3} = [2;0]\)}{\(X = [0;0]\)}{\(X_{1} = [0;0]\),
\(X_{2} = [-1;0]\)}{\(X_{1} = [0;0]\),
\(X_{2} = [1;0]\),
\(X_{3} = [-2;0]\)}{}{}}
\LANGpl{
\Question{
Wyznacz wszystkie punkty przecięcia wykresu podanej funkcji z osią 
\(x\) układu współrzędnych.
\[ 
 f\colon y = x^{3} - x^{2} - 2x
\]}
{\(X_{1} = [0;0]\),
\(X_{2} = [-1;0]\),
\(X_{3} = [2;0]\)}{\(X = [0;0]\)}{\(X_{1} = [0;0]\),
\(X_{2} = [-1;0]\)}{\(X_{1} = [0;0]\),
\(X_{2} = [1;0]\),
\(X_{3} = [-2;0]\)}{}{}}
\LANGcs{
\Question{
Určete všechny společné body osy
\(x\) a grafu
funkce \(f\colon y = x^{3} - x^{2} - 2x\).}
{\(X_{1} = [0;0]\),
\(X_{2} = [-1;0]\),
\(X_{3} = [2;0]\)}{\(X = [0;0]\)}{\(X_{1} = [0;0]\),
\(X_{2} = [-1;0]\)}{\(X_{1} = [0;0]\),
\(X_{2} = [1;0]\),
\(X_{3} = [-2;0]\)}{}{}}
\LANGsk{
\Question{
Určte všetky priesečníky  grafu danej funkcie s osou
\(x\).
\[ 
 f\colon y = x^{3} - x^{2} - 2x
\]}
{\(X_{1} = [0;0]\),
\(X_{2} = [-1;0]\),
\(X_{3} = [2;0]\)}{\(X = [0;0]\)}{\(X_{1} = [0;0]\),
\(X_{2} = [-1;0]\)}{\(X_{1} = [0;0]\),
\(X_{2} = [1;0]\),
\(X_{3} = [-2;0]\)}{}{}}
\LANGes{
\Question{
Determina todas las intersecciones del eje \(x\)  con la gráfica de la siguiente función.
\[ 
 f(x) = x^{3} - x^{2} - 2x
\]}
{\(X_{1} = [0;0]\),
\(X_{2} = [-1;0]\),
\(X_{3} = [2;0]\)}{\(X = [0;0]\)}{\(X_{1} = [0;0]\),
\(X_{2} = [-1;0]\)}{\(X_{1} = [0;0]\),
\(X_{2} = [1;0]\),
\(X_{3} = [-2;0]\)}{}{}}
\End

\Data\ID{1003164101}
\Updated{2020-07-15 03:07:21}
\Level{B}
\SubArea{Power and radical functions}
\LANGen{
\Question{
Let \( f(x)=\sqrt x \). Identify which of the following statements is false.}
{\( f(0.144)=0.12 \)}{\( f(48400)=220 \)}{\( f\!\left(\frac15\right)=\frac{\sqrt5}5 \)}{\( f(588)=14\cdot\sqrt3 \)}{}{}}
\LANGpl{
\Question{
Niech  \( f(x)=\sqrt x \). Które z podanych stwierdzeń jest fałszywe?}
{\( f(0{,}144)=0{,}12 \)}{\( f(48400)=220 \)}{\( f\!\left(\frac15\right)=\frac{\sqrt5}5 \)}{\( f(588)=14\cdot\sqrt3 \)}{}{}}
\LANGcs{
\Question{
Funkce  \( f \) je dána předpisem \( f(x)=\sqrt x \). Vyberte nepravdivý  výrok.}
{\( f(0{,}144)=0{,}12 \)}{\( f(48400)=220 \)}{\( f\!\left(\frac15\right)=\frac{\sqrt5}5 \)}{\( f(588)=14\cdot\sqrt3 \)}{}{}}
\LANGsk{
\Question{
Funkcia \(f\) je daná predpisom \( f(x)=\sqrt x \). Vyberte nepravdivý výrok.}
{\( f(0{,}144)=0{,}12 \)}{\( f(48400)=220 \)}{\( f\!\left(\frac15\right)=\frac{\sqrt5}5 \)}{\( f(588)=14\cdot\sqrt3 \)}{}{}}
\LANGes{
\Question{
Sea \( f(x)=\sqrt x \). Identifica cuál de las siguientes proposiciones lógicas es falsa.}
{\( f(0.144)=0.12 \)}{\( f(48400)=220 \)}{\( f\!\left(\frac15\right)=\frac{\sqrt5}5 \)}{\( f(588)=14\cdot\sqrt3 \)}{}{}}
\End

\Data\ID{1003164102}
\Updated{2020-07-15 03:07:21}
\Level{B}
\SubArea{Power and radical functions}
\LANGen{
\Question{
Let \( f(x)=\sqrt[3]x \). Identify which of the following statements is true.}
{\( 2\cdot f(5)=\sqrt[3]{40} \)}{\( f(0.27)=0.3 \)}{\( f\!\left(\frac13\right)=\frac{\sqrt[3]3}3 \)}{\( f(504)=6\cdot\sqrt[3]7 \)}{}{}}
\LANGpl{
\Question{
Niech \( f(x)=\sqrt[3]x \). Które z poniższych stwierdzeń jest prawdziwe?}
{\( 2\cdot f(5)=\sqrt[3]{40} \)}{\( f(0{,}27)=0{,}3 \)}{\( f\!\left(\frac13\right)=\frac{\sqrt[3]3}3 \)}{\( f(504)=6\cdot\sqrt[3]7 \)}{}{}}
\LANGcs{
\Question{
Funkce \( f \) je dána předpisem \( f(x)=\sqrt[3]x \). Vyberte pravdivý výrok.}
{\( 2\cdot f(5)=\sqrt[3]{40} \)}{\( f(0{,}27)=0{,}3 \)}{\( f\!\left(\frac13\right)=\frac{\sqrt[3]3}3 \)}{\( f(504)=6\cdot\sqrt[3]7 \)}{}{}}
\LANGsk{
\Question{
Funkcia \(f\) je daná predpisom \( f(x)=\sqrt[3]x \). Vyberte pravdivý výrok.}
{\( 2\cdot f(5)=\sqrt[3]{40} \)}{\( f(0{,}27)=0{,}3 \)}{\( f\!\left(\frac13\right)=\frac{\sqrt[3]3}3 \)}{\( f(504)=6\cdot\sqrt[3]7 \)}{}{}}
\LANGes{
\Question{
Sea \( f(x)=\sqrt[3]x \). Identifica cuál de las siguientes proposiciones lógicas es verdadera.}
{\( 2\cdot f(5)=\sqrt[3]{40} \)}{\( f(0.27)=0.3 \)}{\( f\!\left(\frac13\right)=\frac{\sqrt[3]3}3 \)}{\( f(504)=6\cdot\sqrt[3]7 \)}{}{}}
\End

\Data\ID{1003164103}
\Updated{2020-07-15 03:07:21}
\Level{B}
\SubArea{Power and radical functions}
\LANGen{
\Question{
Let \( f(x)=\sqrt[4]x \). Identify which of the following statements is false.}
{\( f\!\left(\sqrt5\right)=\sqrt[6]5 \)}{\( f\!\left(\sqrt[3]{81}\right)=\sqrt[3]3 \)}{\( f(49)=\sqrt7 \)}{\( f\!\left(\frac{25}{256}\right)=\frac{\sqrt5}4 \)}{}{}}
\LANGpl{
\Question{
Niech \( f(x)=\sqrt[4]x \). Które z poniższych stwierdzeń jest   fałszywe?}
{\( f\!\left(\sqrt5\right)=\sqrt[6]5 \)}{\( f\!\left(\sqrt[3]{81}\right)=\sqrt[3]3 \)}{\( f(49)=\sqrt7 \)}{\( f\!\left(\frac{25}{256}\right)=\frac{\sqrt5}4 \)}{}{}}
\LANGcs{
\Question{
Funkce \( f \) je dána předpisem \( f(x)=\sqrt[4]x \). Vyberte nepravdivý výrok.}
{\( f\!\left(\sqrt5\right)=\sqrt[6]5 \)}{\( f\!\left(\sqrt[3]{81}\right)=\sqrt[3]3 \)}{\( f(49)=\sqrt7 \)}{\( f\!\left(\frac{25}{256}\right)=\frac{\sqrt5}4 \)}{}{}}
\LANGsk{
\Question{
Funkcia \(f\) je daná predpisom \( f(x)=\sqrt[4]x \). Vyberte nepravdivý výrok.}
{\( f\!\left(\sqrt5\right)=\sqrt[6]5 \)}{\( f\!\left(\sqrt[3]{81}\right)=\sqrt[3]3 \)}{\( f(49)=\sqrt7 \)}{\( f\!\left(\frac{25}{256}\right)=\frac{\sqrt5}4 \)}{}{}}
\LANGes{
\Question{
Sea \( f(x)=\sqrt[4]x \). Identifica cuál de las siguientes proposiciones lógicas es falsa.}
{\( f\!\left(\sqrt5\right)=\sqrt[6]5 \)}{\( f\!\left(\sqrt[3]{81}\right)=\sqrt[3]3 \)}{\( f(49)=\sqrt7 \)}{\( f\!\left(\frac{25}{256}\right)=\frac{\sqrt5}4 \)}{}{}}
\End

\Data\ID{1003164104}
\Updated{2020-07-15 03:07:21}
\Level{B}
\SubArea{Power and radical functions}
\LANGen{
\Question{
Let \( f(x)=\sqrt x \)  and \( g(x)=\sqrt[3]x \). Identify which of the following statements is true.}
{\( \frac{g(48)}{g(6)} =2 \)}{\( f(15)+g(17)=2 \)}{\( f(2)\cdot g(16)=2 \)}{\( f(11)-f(7)=2 \)}{}{}}
\LANGpl{
\Question{
Niech \( f(x)=\sqrt x \)  i \( g(x)=\sqrt[3]x \). Które z poniższych stwierdzeń jest prawdziwe?}
{\( \frac{g(48)}{g(6)} =2 \)}{\( f(15)+g(17)=2 \)}{\( f(2)\cdot g(16)=2 \)}{\( f(11)-f(7)=2 \)}{}{}}
\LANGcs{
\Question{
Funkce \( f \) a \( g \) jsou dány předpisy  \( f(x)=\sqrt x \)  a \( g(x)=\sqrt[3]x \). Vyberte pravdivý výrok.}
{\( \frac{g(48)}{g(6)} =2 \)}{\( f(15)+g(17)=2 \)}{\( f(2)\cdot g(16)=2 \)}{\( f(11)-f(7)=2 \)}{}{}}
\LANGsk{
\Question{
Funkcie \(f\) a \(g\) sú dané predpismi \( f(x)=\sqrt x \)  a \( g(x)=\sqrt[3]x \). Vyberte pravdivý výrok.}
{\( \frac{g(48)}{g(6)} =2 \)}{\( f(15)+g(17)=2 \)}{\( f(2)\cdot g(16)=2 \)}{\( f(11)-f(7)=2 \)}{}{}}
\LANGes{
\Question{
Sea \( f(x)=\sqrt x \)  and \( g(x)=\sqrt[3]x \). Identifica cuál de las siguientes proposiciones lógicas es verdadera.}
{\( \frac{g(48)}{g(6)} =2 \)}{\( f(15)+g(17)=2 \)}{\( f(2)\cdot g(16)=2 \)}{\( f(11)-f(7)=2 \)}{}{}}
\End

\Data\ID{1003164105}
\Updated{2021-04-12 08:04:52}
\Level{B}
\SubArea{Power and radical functions}
\LANGen{
\Question{
Let \( f(x)=\sqrt{x^2} \) and \( g(x)=x \). Identify which of the following statements is true.}
{\( f(-3)=3 \)}{\( f(-5)=-5 \)}{\(\forall x \in \mathbb{R}:  f(x)=g(x) \)}{\(\forall x \in \mathbb{R}:  f(x)=-g(x) \)}{}{}}
\LANGpl{
\Question{
Niech \( f(x)=\sqrt{x^2} \) i \( g(x)=x \). Które z poniższych stwierdzeń jest prawdziwe?}
{\( f(-3)=3 \)}{\( f(-5)=-5 \)}{\(\forall x \in \mathbb{R}:  f(x)=g(x) \)}{\(\forall x \in \mathbb{R}:  f(x)=-g(x) \)}{}{}}
\LANGcs{
\Question{
Funkce \( f \) a \( g \) jsou dány předpisy \( f(x)=\sqrt{x^2} \) a \( g(x)=x \). Vyberte pravdivý výrok.}
{\( f(-3)=3 \)}{\( f(-5)=-5 \)}{\(\forall x \in \mathbb{R}: f(x)=g(x) \)}{\(\forall x \in \mathbb{R}: f(x)=-g(x) \)}{}{}}
\LANGsk{
\Question{
Funkcie \(f\) a \(g\) sú dané predpismi \( f(x)=\sqrt{x^2} \) a \( g(x)=x \). Vyberte pravdivý výrok.}
{\( f(-3)=3 \)}{\( f(-5)=-5 \)}{\(\forall x \in \mathbb{R}:  f(x)=g(x) \)}{\(\forall x \in \mathbb{R}:  f(x)=-g(x) \)}{}{}}
\LANGes{
\Question{
Sean \( f(x)=\sqrt{x^2} \) y \( g(x)=x \). Identifica cuál de las siguientes proposiciones lógicas es verdadera.}
{\( f(-3)=3 \)}{\( f(-5)=-5 \)}{\(\forall x \in \mathbb{R}:  f(x)=g(x) \)}{\(\forall x \in \mathbb{R}:  f(x)=-g(x) \)}{}{}}
\End

\Data\ID{1003199901}
\Updated{2020-07-15 03:07:21}
\Level{B}
\SubArea{Power and radical functions}
\LANGen{
\Question{
Let \( f(x)=x^{\frac12} \). Identify which of the following statements is false.}
{\( f(0.25)=16 \)}{\( f(0.0121)=0.11 \)}{\( f(\frac12)=\frac{\sqrt2}2 \)}{\( f(338)=13\sqrt2 \)}{}{}}
\LANGpl{
\Question{
Niech \( f(x)=x^{\frac12} \). Które z poniższych stwierdzeń jest fałszywe?}
{\( f(0{,}25)=16 \)}{\( f(0{,}0121)=0{,}11 \)}{\( f(\frac12)=\frac{\sqrt2}2 \)}{\( f(338)=13\sqrt2 \)}{}{}}
\LANGcs{
\Question{
Funkce \( f \) je dána předpisem \( f(x)=x^{\frac12} \). Vyberte nepravdivý výrok.}
{\( f(0{,}25)=16 \)}{\( f(0{,}0121)=0{,}11 \)}{\( f(\frac12)=\frac{\sqrt2}2 \)}{\( f(338)=13\sqrt2 \)}{}{}}
\LANGsk{
\Question{
Funkcia \(f\) je daná predpisom \( f(x)=x^{\frac12} \). Vyberte nepravdivý výrok.}
{\( f(0{,}25)=16 \)}{\( f(0{,}0121)=0{,}11 \)}{\( f(\frac12)=\frac{\sqrt2}2 \)}{\( f(338)=13\sqrt2 \)}{}{}}
\LANGes{
\Question{
Sea \( f(x)=x^{\frac12} \). Identifica cuál de las siguientes proposiciones lógicas es falsa.}
{\( f(0.25)=16 \)}{\( f(0.0121)=0.11 \)}{\( f(\frac12)=\frac{\sqrt2}2 \)}{\( f(338)=13\sqrt2 \)}{}{}}
\End

\Data\ID{1003199902}
\Updated{2020-07-15 03:07:21}
\Level{B}
\SubArea{Power and radical functions}
\LANGen{
\Question{
Let \( f(x)=x^{\frac23} \). Identify which of the following statements is true.}
{\( f(27)=9 \)}{\( f(16)=64 \)}{\( f(\frac18)=4 \)}{\( f(0.0125)=0.25 \)}{}{}}
\LANGpl{
\Question{
Niech \( f(x)=x^{\frac23} \). Które z poniższych stwierdzeń jest prawdziwe?}
{\( f(27)=9 \)}{\( f(16)=64 \)}{\( f(\frac18)=4 \)}{\( f(0{,}0125)=0{,}25 \)}{}{}}
\LANGcs{
\Question{
Funkce \( f \) je dána předpisem \( f(x)=x^{\frac23} \). Vyberte pravdivý výrok.}
{\( f(27)=9 \)}{\( f(16)=64 \)}{\( f(\frac18)=4 \)}{\( f(0{,}0125)=0{,}25 \)}{}{}}
\LANGsk{
\Question{
Funkcia \(f\) je daná predpisom \( f(x)=x^{\frac23} \). Vyberte pravdivý výrok.}
{\( f(27)=9 \)}{\( f(16)=64 \)}{\( f(\frac18)=4 \)}{\( f(0{,}0125)=0{,}25 \)}{}{}}
\LANGes{
\Question{
Sea \( f(x)=x^{\frac23} \). Identifica cuál de las siguientes proposiciones lógicas es verdadera.}
{\( f(27)=9 \)}{\( f(16)=64 \)}{\( f(\frac18)=4 \)}{\( f(0.0125)=0.25 \)}{}{}}
\End

\Data\ID{1003199903}
\Updated{2020-07-15 03:07:21}
\Level{B}
\SubArea{Power and radical functions}
\LANGen{
\Question{
Let \( f(x)=x^{-\frac32} \). Identify which of the following statements is false.}
{\( f(8)=0.25 \)}{\( f(100)=0.001 \)}{\( f(\frac1{16})=64 \)}{\( f(\frac12)=2\sqrt2 \)}{}{}}
\LANGpl{
\Question{
Niech \( f(x)=x^{-\frac32} \). Które z poniższych stwierdzeń jest fałszywe?}
{\( f(8)=0{,}25 \)}{\( f(100)=0{,}001 \)}{\( f(\frac1{16})=64 \)}{\( f(\frac12)=2\sqrt2 \)}{}{}}
\LANGcs{
\Question{
Funkce  \( f \) je dána předpisem \( f(x)=x^{-\frac32} \). Vyberte nepravdivý výrok.}
{\( f(8)=0{,}25 \)}{\( f(100)=0{,}001 \)}{\( f(\frac1{16})=64 \)}{\( f(\frac12)=2\sqrt2 \)}{}{}}
\LANGsk{
\Question{
Funkcia \(f\) je daná predpisom \( f(x)=x^{-\frac32} \). Vyberte nepravdivý výrok.}
{\( f(8)=0{,}25 \)}{\( f(100)=0{,}001 \)}{\( f(\frac1{16})=64 \)}{\( f(\frac12)=2\sqrt2 \)}{}{}}
\LANGes{
\Question{
Sea \( f(x)=x^{-\frac32} \). Identifica cuál de las siguientes proposiciones lógicas es falsa.}
{\( f(8)=0.25 \)}{\( f(100)=0.001 \)}{\( f(\frac1{16})=64 \)}{\( f(\frac12)=2\sqrt2 \)}{}{}}
\End

\Data\ID{1003199904}
\Updated{2020-07-15 03:07:21}
\Level{B}
\SubArea{Power and radical functions}
\LANGen{
\Question{
Let \( f(x)=x^{\frac12} \) and \( g(x)=x^{\frac32} \). Identify which of the following statements is true.}
{\( \frac{f(2)}{g(2)} =0.5 \)}{\( f(2)+g(2)=4 \)}{\( f(4)\cdot g(4)=8 \)}{\( \bigl(f(2)\bigr)^{\frac32}=4 \)}{}{}}
\LANGpl{
\Question{
Niech \( f(x)=x^{\frac12} \) i \( g(x)=x^{\frac32} \). Które z poniższych stwierdzeń jest prawdziwe?}
{\( \frac{f(2)}{g(2)} =0{,}5 \)}{\( f(2)+g(2)=4 \)}{\( f(4)\cdot g(4)=8 \)}{\( \bigl(f(2)\bigr)^{\frac32}=4 \)}{}{}}
\LANGcs{
\Question{
Funkce \( f \) a \( g \) jsou dány předpisy \( f(x)=x^{\frac12} \); \( g(x)=x^{\frac32} \). Vyberte pravdivý výrok.}
{\( \frac{f(2)}{g(2)} =0{,}5 \)}{\( f(2)+g(2)=4 \)}{\( f(4)\cdot g(4)=8 \)}{\( \bigl(f(2)\bigr)^{\frac32}=4 \)}{}{}}
\LANGsk{
\Question{
Funkcie \(f\) a \(g\) sú dané predpismi \( f(x)=x^{\frac12} \) a \( g(x)=x^{\frac32} \). Vyberte pravdivý výrok.}
{\( \frac{f(2)}{g(2)} =0{,}5 \)}{\( f(2)+g(2)=4 \)}{\( f(4)\cdot g(4)=8 \)}{\( \bigl(f(2)\bigr)^{\frac32}=4 \)}{}{}}
\LANGes{
\Question{
Sean \( f(x)=x^{\frac12} \) y \( g(x)=x^{\frac32} \). Identifica cuál de las siguientes proposiciones lógicas es verdadera.}
{\( \frac{f(2)}{g(2)} =0.5 \)}{\( f(2)+g(2)=4 \)}{\( f(4)\cdot g(4)=8 \)}{\( \bigl(f(2)\bigr)^{\frac32}=4 \)}{}{}}
\End

\Data\ID{1003199905}
\Updated{2020-07-15 03:07:21}
\Level{B}
\SubArea{Power and radical functions}
\LANGen{
\Question{
Let \( f(x)=\left(\frac{x^2}{\sqrt[3]x}\right)^{-0.5} \). Identify which of the following statements is true.}
{\( f(x)=x^{-\frac56} \)}{\( f(x)=x^{\frac13} \)}{\( f(x)=x^{-\frac16} \)}{\( f(x)=x^{\frac76} \)}{}{}}
\LANGpl{
\Question{
Niech \( f(x)=\left(\frac{x^2}{\sqrt[3]x}\right)^{-0{,}5} \). Które z poniższych stwierdzeń jest prawdziwe?}
{\( f(x)=x^{-\frac56} \)}{\( f(x)=x^{\frac13} \)}{\( f(x)=x^{-\frac16} \)}{\( f(x)=x^{\frac76} \)}{}{}}
\LANGcs{
\Question{
Funkce \( f \) je dána předpisem \( f(x)=\left(\frac{x^2}{\sqrt[3]x}\right)^{-0{,}5} \). Vyberte pravdivý výrok.}
{\( f(x)=x^{-\frac56} \)}{\( f(x)=x^{\frac13} \)}{\( f(x)=x^{-\frac16} \)}{\( f(x)=x^{\frac76} \)}{}{}}
\LANGsk{
\Question{
Funkcia \(f\) je daná predpisom \( f(x)=\left(\frac{x^2}{\sqrt[3]x}\right)^{-0{,}5} \). Vyberte pravdivý výrok.}
{\( f(x)=x^{-\frac56} \)}{\( f(x)=x^{\frac13} \)}{\( f(x)=x^{-\frac16} \)}{\( f(x)=x^{\frac76} \)}{}{}}
\LANGes{
\Question{
Sea \( f(x)=\left(\frac{x^2}{\sqrt[3]x}\right)^{-0.5} \). Identifica cuál de las siguientes proposiciones lógicas es verdadera.}
{\( f(x)=x^{-\frac56} \)}{\( f(x)=x^{\frac13} \)}{\( f(x)=x^{-\frac16} \)}{\( f(x)=x^{\frac76} \)}{}{}}
\End

\Data\ID{1103186201}
\Updated{2020-07-15 03:07:10}
\Level{B}
\SubArea{Power and radical functions}
\LANGen{
\Question{
Identify which of the graphs represents a part of the function \( f(x)=\sqrt{-x} \).}
{}{}{}{}{}{}}
\LANGpl{
\Question{
Który z wykresów przedstawia część wykresu funkcji \( f(x)=\sqrt{-x} \)?}
{}{}{}{}{}{}}
\LANGcs{
\Question{
Na kterém obrázku je část grafu funkce \( f: y=\sqrt{-x} \) ?}
{}{}{}{}{}{}}
\LANGsk{
\Question{
Na ktorom z obrázkov sa nachádza časť grafu funkcie \( f(x)=\sqrt{-x} \)?}
{}{}{}{}{}{}}
\LANGes{
\Question{
Identifica cuál de las gráficas representa una parte de la función \( f(x)=\sqrt{-x} \).}
{}{}{}{}{}{}}

\IMGanswer{1103186201a_0.pdf}
\IMGanswer{1103186201b_0.pdf}
\IMGanswer{1103186201c_0.pdf}
\IMGanswer{1103186201d_0.pdf}\End

\Data\ID{9000021708}
\Updated{2021-04-07 09:04:18}
\Level{B}
\SubArea{Power and radical functions}
\LANGen{
\Question{
In the following list identify a function with a domain
\(\left (-\infty ; \frac{1} 
{2}\right )\).}
{\(f(x)= \sqrt{ \frac{10}
{2-4x}}\)}{\(f(x)= \sqrt{\frac{2-4x}
 {10}} \)}{\(f(x) = \sqrt{2 - 4x}\)}{\(f(x) = \sqrt{\frac{2-4x}
 {3x}} \)}{}{}}
\LANGpl{
\Question{
Dziedziną, której z poniższych funkcji jest przedział
\(\left (-\infty ; \frac{1} 
{2}\right )\)?}
{\(f\colon y = \sqrt{ \frac{10}
{2-4x}}\)}{\(f\colon y = \sqrt{\frac{2-4x}
 {10}} \)}{\(f\colon y = \sqrt{2 - 4x}\)}{\(f\colon y = \sqrt{\frac{2-4x}
 {3x}} \)}{}{}}
\LANGcs{
\Question{
Interval \(\left (-\infty ; \frac{1} 
{2}\right )\) 
je definičním oborem funkce:}
{\(f\colon y = \sqrt{ \frac{10}
{2-4x}}\)}{\(f\colon y = \sqrt{\frac{2-4x}
 {10}} \)}{\(f\colon y = \sqrt{2 - 4x}\)}{\(f\colon y = \sqrt{\frac{2-4x}
 {3x}} \)}{}{}}
\LANGsk{
\Question{
Určte funkciu, ktorej definičný obor je interval \(\left (-\infty ; \frac{1} 
{2}\right )\).}
{\(f\colon y = \sqrt{ \frac{10}
{2-4x}}\)}{\(f\colon y = \sqrt{\frac{2-4x}
 {10}} \)}{\(f\colon y = \sqrt{2 - 4x}\)}{\(f\colon y = \sqrt{\frac{2-4x}
 {3x}} \)}{}{}}
\LANGes{
\Question{
En la siguiente lista, identifica una función cuyo dominio sea
\(\left (-\infty ; \frac{1} 
{2}\right )\).}
{\(f(x)= \sqrt{ \frac{10}
{2-4x}}\)}{\(f(x)= \sqrt{\frac{2-4x}
 {10}} \)}{\(f(x) = \sqrt{2 - 4x}\)}{\(f(x) = \sqrt{\frac{2-4x}
 {3x}} \)}{}{}}
\End

\Data\ID{9000025804}
\Updated{2020-07-15 03:07:34}
\Level{B}
\SubArea{Power and radical functions}
\LANGen{
\Question{
In the following list identify a true statement on the function
\(f\).
\[ 
 f(x) = (x + 1)(x + 2)(x - 3)
\]}
{The function \(f\) 
is positive on \(I_{1} = (-2;-1)\) 
and \(I_{2} = (3;\infty )\).}{The function \(f\) 
is an increasing function (in its whole domain).}{The function is decreasing only on
\(I = (-1;3)\).}{The function is decreasing on \(I_{1} = (-\infty ;-2)\) 
and \(I_{2} = (3;\infty )\).}{}{}}
\LANGpl{
\Question{
Które ze stwierdzeń dotyczących funkcji \(f\)  jest prawdziwe?
\[ 
 f\colon y = (x + 1)(x + 2)(x - 3)
\]}
{Funkcja  \(f\) 
przyjmuje wartości dodatnie dla \(I_{1} = (-2;-1)\) 
i \(I_{2} = (3;\infty )\).}{Funkcja \(f\) 
jest funkcją rosnącą (w całej dziedzinie).}{Funkcja jest malejąca tylko dla
\(I = (-1;3)\).}{Funkcja jest malejąca dla \(I_{1} = (-\infty ;-2)\) 
i \(I_{2} = (3;\infty )\).}{}{}}
\LANGcs{
\Question{
Který z následujících výroků o funkci
\(f\colon y = (x + 1)(x + 2)(x - 3)\) je
pravdivý?}
{Funkce nabývá kladných hodnot právě ve dvou intervalech
\(I_{1} = (-2;-1)\) a
\(I_{2} = (3;\infty )\).}{Funkce je rostoucí v celém \(D(f)\).}{Funkce je klesající pouze v intervalu
\(I = (-1;3)\).}{Funkce je klesající právě ve dvou intervalech
\(I_{1} = (-\infty ;-2)\) a
\(I_{2} = (3;\infty )\).}{}{}}
\LANGsk{
\Question{
Ktorý z následujúcich výrokov o funkcii
\(f\colon y = (x + 1)(x + 2)(x - 3)\) je
pravdivý?}
{Funkcia nadobúda kladné hodnoty práve na dvoch intervaloch
\(I_{1} = (-2;-1)\) a
\(I_{2} = (3;\infty )\).}{Funkcia je rastúca na celom \(D(f)\).}{Funkcia je klesajúca len na intervale
\(I = (-1;3)\).}{Funkcia je klesajúca práve na dvoch intervaloch
\(I_{1} = (-\infty ;-2)\) a
\(I_{2} = (3;\infty )\).}{}{}}
\LANGes{
\Question{
En la siguiente lista, identifica una proposición lógica verdadera sobre la función
\(f\).
\[ 
 f(x) = (x + 1)(x + 2)(x - 3)
\]}
{La función \(f\) 
es positiva en \(I_{1} = (-2;-1)\) y  \(I_{2} = (3;\infty )\).}{La función \(f\) 
es una función creciente (en todo su dominio).}{La función está disminuyendo solo en
\(I = (-1;3)\).}{La función  está disminuyendo solo en \(I_{1} = (-\infty ;-2)\) 
y \(I_{2} = (3;\infty )\).}{}{}}
\End

\Data\ID{1003101101}
\Updated{2021-04-07 09:04:27}
\Level{C}
\SubArea{Power and radical functions}
\LANGen{
\Question{
Find the true statement about the function  \( f(x)=\left|x^3+1\right| \).}
{The function \( f \) has the minimum at \( x=-1 \).}{The function \( f \) has the minimum at \( x=0 \).}{The function \( f \) has the minimum at \( x=1 \).}{The function \( f \) has no minimum.}{}{}}
\LANGpl{
\Question{
Które z poniższych stwierdzeń dotyczących funkcji \( f(x)=\left|x^3+1\right| \) jest prawdziwe?}
{Funkcja \( f \) osiąga minimum w \( x=-1 \).}{Funkcja \( f \) osiąga minimum w \( x=0 \).}{Funkcja \( f \) osiąga minimum w \( x=1 \).}{Funkcja \( f \) nie ma minimum.}{}{}}
\LANGcs{
\Question{
Funkce \( f \) je dána předpisem  \( f(x)=\left|x^3+1\right| \). Vyberte pravdivý výrok.}
{Funkce \( f \) má minimum v bodě \( x=-1 \).}{Funkce \( f \) má minimum v bodě \( x=0 \).}{Funkce \( f \) má minimum v bodě \( x=1 \).}{Funkce \( f \) nemá minimum.}{}{}}
\LANGsk{
\Question{
Daná je funkcia predpisom \( f(x)=\left|x^3+1\right| \). Vyberte pravdivé tvrdenie o funkcii \( f \).}
{Funkcia \( f \) má minimum v bode \( x=-1 \).}{Funkcia \( f \) má minimum v bode \( x=0 \).}{Funkcia \( f \) má minimum v bode \( x=1 \).}{Funkcia \( f \) nemá minimum.}{}{}}
\LANGes{
\Question{
Determina la proposición lógica verdadera sobre la función  \( f(x)=\left|x^3+1\right| \).}
{La función \( f \) tiene el mínimo en \( x=-1 \).}{La función \( f \) tiene el mínimo en \( x=0 \).}{La función \( f \) tiene el mínimo en \( x=1 \).}{La función  \( f \) no tiene mínimo.}{}{}}
\End

\Data\ID{1003159101}
\Updated{2021-06-30 11:06:38}
\Level{C}
\SubArea{Power and radical functions}
\LANGen{
\Question{
Choose the function that describes the dependence of a cube volume \( V \) on the length of its space diagonal \( u \).}
{\( V=\frac{\sqrt3\cdot u^3}9;\ u\in(0;\infty) \)}{\( V=u^3;\ u\in(0;\infty) \)}{\( V=\frac{u^3}3;\ u\in(0;\infty) \)}{\( V=27u^3;\ u\in(0;\infty) \)}{}{}}
\LANGpl{
\Question{
Wybierz funkcję opisującą zależność pomiędzy objętością sześcianu \( V \) a długością  jego przekątnej \( u \).}
{\( V=\frac{\sqrt3\cdot u^3}9;\ u\in(0;\infty) \)}{\( V=u^3;\ u\in(0;\infty) \)}{\( V=\frac{u^3}3;\ u\in(0;\infty) \)}{\( V=27u^3;\ u\in(0;\infty) \)}{}{}}
\LANGcs{
\Question{
Vyberte funkci, která vyjadřuje závislost objemu  krychle \( V \) na délce její tělesové uhlopříčky \( u \).}
{\( V=\frac{\sqrt3\cdot u^3}9;\ u\in(0;\infty) \)}{\( V=u^3;\ u\in(0;\infty) \)}{\( V=\frac{u^3}3;\ u\in(0;\infty) \)}{\( V=27u^3;\ u\in(0;\infty) \)}{}{}}
\LANGsk{
\Question{
Vyberte funkciu, ktorá vyjadruje závislosť objemu kocky \( V \) na dĺžke telesovej uhlopriečky \( u \).}
{\( V=\frac{\sqrt3\cdot u^3}9;\ u\in(0;\infty) \)}{\( V=u^3;\ u\in(0;\infty) \)}{\( V=\frac{u^3}3;\ u\in(0;\infty) \)}{\( V=27u^3;\ u\in(0;\infty) \)}{}{}}
\LANGes{
\Question{
Elige la función que describe la dependencia del volumen de un cubo  \( V \) respecto a la longitud  \( u \) de su diagonal principal.}
{\( V=\frac{\sqrt3\cdot u^3}9;\ u\in(0;\infty) \)}{\( V=u^3;\ u\in(0;\infty) \)}{\( V=\frac{u^3}3;\ u\in(0;\infty) \)}{\( V=27u^3;\ u\in(0;\infty) \)}{}{}}
\End

\Data\ID{1003159201}
\Updated{2021-08-23 02:08:03}
\Level{C}
\SubArea{Power and radical functions}
\LANGen{
\Question{
The 3D printer prints a solid \( 5 \) centimetre cube in \( 2 \) hours. The printer can print the cube with the maximum edge length of \( 20\,\mathrm{cm} \). Suppose the printing time is directly proportional to the cube volume. Choose the function that describes the dependence of the number \( n \) of cubes printed in \( 1 \) day on the printed cube edge length \( a \), which is specified in centimetres. Neglect time needed for using the printer.}
{\( n=1500a^{-3};\ a\in(0;20] \)}{\( n=60a^{-1};\ a\in(0;20] \)}{\( n=300a^{-2};\ a\in(0;20] \)}{\( n=2.4a;\ a\in(0;20] \)}{}{}}
\LANGpl{
\Question{
Drukarka  3D drukuje  \( 5 \) - centymetrowy sześcian  w \( 2 \) godziny. Drukarka może wydrukować sześcian o maksymalnej długości krawędzi \( 20\,\mathrm{cm} \). załóżmy, że czas wydruku jest wprost proporcjonalny do objętości sześcianu. Wybierz funkcję, która opisuje zależność liczby wydrukowanych sześcianów \( n \) w ciągu  \( 1 \) dnia od długości krawędzi wydrukowanego sześcianu \( a \), określonej w centymetrach. Nie bierz pod uwagę czasu drukowania.}
{\( n=1500a^{-3};\ a\in(0;20\rangle \)}{\( n=60a^{-1};\ a\in(0;20\rangle \)}{\( n=300a^{-2};\ a\in(0;20\rangle \)}{\( n=2{,}4a;\ a\in(0;20\rangle \)}{}{}}
\LANGcs{
\Question{
3D tiskárna vytiskne plnou krychli o hraně délky \( 5 \,\mathrm{cm} \) za \( 2 \,\mathrm{hodiny} \). Tiskárna dokáže vytisknout krychli o maximální délce hrany \( 20\,\mathrm{cm} \). Předpokládejme, že doba tisku je přímo úměrná objemu krychle. Vyberte předpis funkce, která vyjadřuje závislost počtu krychlí \( n \) vytištěných za \( 1 \,\mathrm{den} \) na délce hrany \( a \) krychle zadané v centimetrech. Dobu potřebnou k obsluze tiskárny zanedbejte.}
{\( n=1500 a^{-3};\ a\in(0;20\rangle \)}{\( n=60 a^{-1};\ a\in(0;20\rangle \)}{\( n=300 a^{-2};\ a\in(0;20\rangle \)}{\( n=2{,}4 a;\ a\in(0;20\rangle \)}{}{}}
\LANGsk{
\Question{
3D tlačiareň vytlačí plnú kocku s hranou dĺžky \( 5 \,\mathrm{cm} \) za \( 2 \,\mathrm{hodiny} \). Tlačiareň dokáže vytlačiť kocku s maximálnou dĺžkou hrany \( 20\,\mathrm{cm} \). Prepokladajme, že čas tlačenia je priamo úmerný objemu kocky. Vyberte predpis funkcie, ktorá vyjadruje závislosť počtu kociek \( n \) vytlačených za \( 1 \) deň na dĺžke hrany \( a \) kocky zadanej v centimetroch. Čas potrebný na obsluhu tlačiarne zanedbajte.}
{\( n=1500 a^{-3};\ a\in(0;20\rangle \)}{\( n=60 a^{-1};\ a\in(0;20\rangle \)}{\( n=300 a^{-2};\ a\in(0;20\rangle \)}{\( n=2{,}4 a;\ a\in(0;20\rangle \)}{}{}}
\LANGes{
\Question{
Una impresora 3D imprime un cubo sólido de  \( 5 \) centímetros en \( 2 \) horas. La impresora puede imprimir un cubo con una longitud máxima de arista de \( 20\,\mathrm{cm} \). Supongamos que el tiempo de la impresión es directamente proporcional al volumen del cubo. Elige la función que describe la dependencia del número \( n \) de cubos impresos en  \( 1 \) día respecto de la longitud de la arista del cubo impreso \( a\),  que se especifica en centímetros. Ignora el tiempo necesario para usar la impresora.}
{\( n=1500a^{-3};\ a\in(0;20] \)}{\( n=60a^{-1};\ a\in(0;20] \)}{\( n=300a^{-2};\ a\in(0;20] \)}{\( n=2.4a;\ a\in(0;20] \)}{}{}}
\End

\Data\ID{1103082706}
\Updated{2020-07-15 03:07:10}
\Level{C}
\SubArea{Power and radical functions}
\LANGen{
\Question{
Identify the graph of the function \( f(x)=\sqrt{|-x|};\ x\in [-9;9] \).}
{}{}{}{}{}{}}
\LANGpl{
\Question{
Wybierz wykres funkcji \( f(x)=\sqrt{|-x|};\ x\in\langle-9;9\rangle \).}
{}{}{}{}{}{}}
\LANGcs{
\Question{
Určete graf funkce \( f(x)=\sqrt{|-x|};\ x\in\langle-9;9\rangle \).}
{}{}{}{}{}{}}
\LANGsk{
\Question{
Ktorý z následujúcich grafov je graf funkcie \( f(x)=\sqrt{|-x|};\ x\in\langle-9;9\rangle \)?}
{}{}{}{}{}{}}
\LANGes{
\Question{
Identifica la gráfica de la función \( f(x)=\sqrt{|-x|};\ x\in [-9;9] \).}
{}{}{}{}{}{}}

\IMGanswer{1103082706a.pdf}
\IMGanswer{1103082706b.pdf}
\IMGanswer{1103082706c.pdf}
\IMGanswer{1103082706d.pdf}\End
